% =============================================================================
%  MARCO TEÓRICO — Caracterización Dieléctrica mediante Sonda Coaxial Abierta
%  Proyecto Final de Ingeniería Electrónica — UTN FRBA
%  Axel Nahum — 2026
%
%  Este archivo contiene únicamente el contenido de la sección teórica.
%  Para compilarlo de forma standalone basta con envolver el bloque en:
%
%    \documentclass[12pt,a4paper]{article}
%    \usepackage[utf8]{inputenc}
%    \usepackage[T1]{fontenc}
%    \usepackage[spanish]{babel}
%    \usepackage{amsmath,amssymb,amsfonts}
%    \usepackage{booktabs}
%    \usepackage{graphicx}
%    \usepackage{hyperref}
%    \begin{document}
%    % =============================================================================
%  MARCO TEÓRICO — Caracterización Dieléctrica mediante Sonda Coaxial Abierta
%  Proyecto Final de Ingeniería Electrónica — UTN FRBA
%  Axel Nahum — 2026
%
%  Este archivo contiene únicamente el contenido de la sección teórica.
%  Para compilarlo de forma standalone basta con envolver el bloque en:
%
%    \documentclass[12pt,a4paper]{article}
%    \usepackage[utf8]{inputenc}
%    \usepackage[T1]{fontenc}
%    \usepackage[spanish]{babel}
%    \usepackage{amsmath,amssymb,amsfonts}
%    \usepackage{booktabs}
%    \usepackage{graphicx}
%    \usepackage{hyperref}
%    \begin{document}
%    % =============================================================================
%  MARCO TEÓRICO — Caracterización Dieléctrica mediante Sonda Coaxial Abierta
%  Proyecto Final de Ingeniería Electrónica — UTN FRBA
%  Axel Nahum — 2026
%
%  Este archivo contiene únicamente el contenido de la sección teórica.
%  Para compilarlo de forma standalone basta con envolver el bloque en:
%
%    \documentclass[12pt,a4paper]{article}
%    \usepackage[utf8]{inputenc}
%    \usepackage[T1]{fontenc}
%    \usepackage[spanish]{babel}
%    \usepackage{amsmath,amssymb,amsfonts}
%    \usepackage{booktabs}
%    \usepackage{graphicx}
%    \usepackage{hyperref}
%    \begin{document}
%    % =============================================================================
%  MARCO TEÓRICO — Caracterización Dieléctrica mediante Sonda Coaxial Abierta
%  Proyecto Final de Ingeniería Electrónica — UTN FRBA
%  Axel Nahum — 2026
%
%  Este archivo contiene únicamente el contenido de la sección teórica.
%  Para compilarlo de forma standalone basta con envolver el bloque en:
%
%    \documentclass[12pt,a4paper]{article}
%    \usepackage[utf8]{inputenc}
%    \usepackage[T1]{fontenc}
%    \usepackage[spanish]{babel}
%    \usepackage{amsmath,amssymb,amsfonts}
%    \usepackage{booktabs}
%    \usepackage{graphicx}
%    \usepackage{hyperref}
%    \begin{document}
%    \input{teoria_caracterizacion}
%    \end{document}
% =============================================================================

\section{Marco Teórico: Caracterización Dieléctrica mediante Sonda Coaxial Abierta}

El presente capítulo desarrolla los fundamentos físicos, matemáticos y metrológicos
del método de sonda coaxial abierta empleado por el sistema NanoVNA-UTN-Toolkit
para la determinación de la permitividad eléctrica compleja de materiales líquidos.
Se presentan los dos algoritmos implementados —el método completo y el simplificado—,
los modelos de Debye utilizados para los líquidos de referencia, y los criterios de
selección de dichos líquidos.

\newpage
\Encabezado

% -----------------------------------------------------------------------------
\subsection{La Permitividad Eléctrica Compleja}
% -----------------------------------------------------------------------------

La permitividad eléctrica $\varepsilon$ es la magnitud que describe la respuesta de
un material a la aplicación de un campo eléctrico externo. En el dominio frecuencial,
para un material dieléctrico con pérdidas, se expresa como un número complejo:

\begin{equation}
    \varepsilon^*(\omega) = \varepsilon'(\omega) - j\,\varepsilon''(\omega),
    \label{eq:permitividad_compleja}
\end{equation}

\noindent donde:
\begin{itemize}
    \item $\varepsilon'$ es la parte real, asociada a la energía almacenada en el campo
          eléctrico (capacitiva). Cuantifica cuánto puede polarizarse el material.
    \item $\varepsilon''$ es la parte imaginaria, siempre no negativa para un medio
          pasivo, asociada a las pérdidas dieléctricas (conversión de energía
          electromagnética en calor).
    \item $\omega = 2\pi f$ es la frecuencia angular.
\end{itemize}

La permitividad relativa $\varepsilon_r$ se define normalizando respecto a la
permitividad del vacío $\varepsilon_0 = 8{,}854 \times 10^{-12}$~F/m:

\begin{equation}
    \varepsilon_r^* = \frac{\varepsilon^*}{\varepsilon_0}
                    = \varepsilon_r' - j\,\varepsilon_r''.
\end{equation}

Una magnitud derivada de uso frecuente es el \emph{factor de disipación} o
\emph{tangente de pérdidas}:

\begin{equation}
    \tan\delta = \frac{\varepsilon_r''}{\varepsilon_r'}.
\end{equation}

A lo largo de este documento se omite el subíndice $r$ salvo aclaración explícita;
todas las magnitudes son relativas.

\newpage
\Encabezado

% -----------------------------------------------------------------------------
\subsection{Modelo de Debye para Líquidos Polares}
\label{sec:debye}
% -----------------------------------------------------------------------------

El modelo de Debye~\cite{debye1929} describe la dependencia en frecuencia de la
permitividad de un dieléctrico polar cuyas moléculas poseen un único tiempo de
relajación $\tau$. Físicamente, $\tau$ cuantifica el tiempo característico de
reorientación dipolar ante un campo oscilante.

\subsubsection{Modelo de Debye simple}

\begin{equation}
    \varepsilon^*(f) = \varepsilon_\infty
                     + \frac{\varepsilon_s - \varepsilon_\infty}{1 + j\,f / f_r},
    \label{eq:debye_simple}
\end{equation}

\noindent donde:
\begin{itemize}
    \item $\varepsilon_s$ es la permitividad estática (límite $f \to 0$).
    \item $\varepsilon_\infty$ es la permitividad a alta frecuencia (límite $f \to \infty$),
          debida a polarizaciones electrónicas e iónicas.
    \item $f_r = 1/(2\pi\tau)$ es la frecuencia de relajación.
\end{itemize}

Este modelo se aplica al agua y al metanol en el presente sistema.
Los coeficientes $\{\varepsilon_s, \varepsilon_\infty, \tau\}$ dependen de la
temperatura y se tabulan a partir de los datos del NPL Report MAT~23
\cite{gregory2012} mediante ajuste polinomial cuadrático en temperatura.

\subsubsection{Modelo de Debye doble}

Algunos líquidos exhiben dos mecanismos de relajación dominantes a frecuencias
distintas (por ejemplo, el isopropanol / IPA). Se modela con dos términos de Debye
superpuestos:

\begin{equation}
    \varepsilon^*(f) = \varepsilon_\infty
                     + \frac{\varepsilon_s - \varepsilon_h}{1 + j\,f / f_{r1}}
                     + \frac{\varepsilon_h - \varepsilon_\infty}{1 + j\,f / f_{r2}},
    \label{eq:debye_doble}
\end{equation}

\noindent donde $\varepsilon_h$ es una permitividad intermedia y $f_{r1} < f_{r2}$
son las dos frecuencias de relajación.

\subsubsection{Modelo de Debye-Gamma}

Para el etanol, el NPL Report MAT~23 propone un modelo de Debye simple con
corrección lineal en frecuencia (término~$\gamma$) que aproxima una segunda
relajación en la parte superior de la banda:

\begin{equation}
    \varepsilon^*(f) = \varepsilon_h
                     + \frac{\varepsilon_s - \varepsilon_h}{1 + j\,f / f_r}
                     - j\,\gamma \left(\frac{f}{\text{GHz}}\right),
    \label{eq:debye_gamma}
\end{equation}

siendo $\gamma$ un coeficiente dimensional en GHz$^{-1}$.

En la Tabla~\ref{tab:modelos_debye} se resume el modelo utilizado para cada
líquido disponible en el sistema.

\begin{table}[h!]
\centering
\caption{Modelos de Debye por líquido de referencia implementados en el sistema.}
\label{tab:modelos_debye}
\begin{tabular}{@{}llcc@{}}
\toprule
\textbf{Líquido} & \textbf{Modelo} & \textbf{Rango T (°C)} & \textbf{Fuente} \\
\midrule
Agua (H$_2$O) & Debye simple       & 0 -- 60  & NPL MAT 23 \cite{gregory2012} \\
IPA (isopropanol) & Debye doble    & 10 -- 30 & NPL MAT 23 \cite{gregory2012} \\
Etanol         & Debye-Gamma        & 10 -- 50 & NPL MAT 23 \cite{gregory2012} \\
Metanol        & Debye simple       & 10 -- 50 & NPL MAT 23 \cite{gregory2012} \\
\bottomrule
\end{tabular}
\end{table}

\newpage
\Encabezado

% -----------------------------------------------------------------------------
\subsection{Método de la Sonda Coaxial Abierta}
\label{sec:sonda_coaxial}
% -----------------------------------------------------------------------------

\subsubsection{Principio físico}

La técnica de sonda coaxial abierta es un método de reflexión no resonante
y no destructivo para la medición de permitividad~\cite{stuchly1980, misra1990}.
Consiste en apoyar o sumergir el extremo abierto de una línea coaxial sobre
el material bajo prueba (MBP) y medir el coeficiente de reflexión~$\Gamma$
(parámetro~$S_{11}$) en un rango de frecuencias.

Cuando el extremo de la línea coaxial se abre a un semiespacio relleno con
el material, la terminación coaxial abierta actúa como una antena de apertura
en campo cercano. La admitancia en la terminación depende fuertemente de las
propiedades dieléctricas del material en contacto:

\begin{equation}
    Y(\varepsilon^*, \omega)
    = j\omega C_0\,\varepsilon^*
    + G_0\,(\varepsilon^*)^{5/2},
    \label{eq:admitancia_modelo}
\end{equation}

\noindent donde:
\begin{itemize}
    \item $C_0$ [F] es la capacitancia de la sonda en el vacío, determinada por la
          geometría del coaxial. Domina el comportamiento a bajas frecuencias.
    \item $G_0$ [S] es el coeficiente de radiación. El término $G_0\,\varepsilon^{5/2}$
          representa las pérdidas por radiación en el campo lejano; domina a altas
          frecuencias y para sondas de mayor apertura (sonda de 21~mm respecto
          a la de 3~mm).
\end{itemize}

La relación entre la admitancia y el coeficiente de reflexión medido es la
transformada de Möbius estándar de la teoría de redes:

\begin{equation}
    \Gamma = \frac{Z_0\,Y - 1}{Z_0\,Y + 1}
    \quad \Longleftrightarrow \quad
    Y = \frac{1}{Z_0}\cdot\frac{1 - \Gamma}{1 + \Gamma},
\end{equation}

\noindent con $Z_0 = 50~\Omega$.

\subsubsection{Historia y contexto del método}

El método de sonda coaxial abierta fue formalizado matemáticamente en la
década de 1980 por Stuchly et al.~\cite{stuchly1980} y Misra~\cite{misra1990},
quienes propusieron el modelo de admitancia de~\eqref{eq:admitancia_modelo}
y demostraron que el exponente $5/2$ emerge del desarrollo en serie del campo
radiado por una apertura circular en un plano conductor infinito.

La calibración con líquidos de referencia, en lugar de estándares mecánicos de
reflexión, fue introducida para eliminar la dependencia del resultado respecto
a la calidad del contacto sonda-muestra y respecto a los errores sistemáticos
del analizador de redes. La versión de dos líquidos que se documenta aquí es
consistente con la referencia UTN-FRBA 2024~\cite{utnfrba2024} en la que se basa
el presente trabajo.

\newpage
\Encabezado

% -----------------------------------------------------------------------------
\subsection{Calibración del Sistema}
\label{sec:calibracion}
% -----------------------------------------------------------------------------

El sistema presenta dos métodos de calibración con distinto número de estándares,
distintos supuestos y distinta complejidad algorítmica.

\subsubsection{Método Completo (4 estándares)}
\label{sec:metodo_completo}

Este método es denominado en la literatura como calibración de 4 estándares
del tipo SOL + 2 líquidos (Open -- Short -- Liquid\,1 -- Liquid\,2).
Los cuatro estándares son:

\begin{enumerate}
    \item \textbf{Circuito abierto (Open)}: la sonda apunta al aire.
          $\varepsilon_{\text{open}} = \varepsilon_{\text{aire}} \approx 1$.
    \item \textbf{Cortocircuito (Short)}: la apertura de la sonda se cubre con
          una lámina metálica conductora.
    \item \textbf{Líquido de referencia~1 (Ref\,1)}: por defecto IPA (isopropanol),
          de permitividad baja-moderada.
    \item \textbf{Líquido de referencia~2 (Ref\,2)}: por defecto agua destilada,
          de permitividad alta.
\end{enumerate}

Las cuatro mediciones $S_{11}$ de calibración y las permitividades conocidas
$\varepsilon_{\text{Ref1}}$, $\varepsilon_{\text{Ref2}}$ (calculadas mediante
los modelos de Debye de la Sección~\ref{sec:debye} a la temperatura de medición)
permiten determinar los parámetros del sistema $\{C_0, G_0\}$ y los tres
coeficientes de error de la caja de errores del analizador (directividad,
adaptación de fuente, seguimiento de reflexión).

La calibración completa absorbe todos los errores sistemáticos del analizador y
del cable de la sonda simultáneamente con los parámetros de la sonda misma.
No se requiere una calibración OSM previa del analizador.

\subsubsection{Método Simplificado (3 estándares)}
\label{sec:metodo_simplificado}

El método simplificado (denominado también \emph{incompleto} en la literatura)
utiliza únicamente tres estándares:

\begin{enumerate}
    \item \textbf{Open}: ídem al método completo.
    \item \textbf{Short}: ídem al método completo.
    \item \textbf{Líquido de referencia único (Ref\,1)}: por defecto agua destilada.
\end{enumerate}

La simplificación consiste en despreciar el término de radiación ($G_0 = 0$),
reduciendo el modelo de admitancia a uno puramente capacitivo:

\begin{equation}
    Y \approx j\omega C_0\,\varepsilon^*.
\end{equation}

Con esta hipótesis, la relación entre $S_{11}$ y $\varepsilon^*$ se vuelve una
transformada de Möbius (bilineal) exacta. Tres puntos conocidos en el plano~$S_{11}$
(aire, cortocircuito, líquido de referencia) determinan unívocamente la transformada,
y la extracción de $\varepsilon^*$ del material desconocido se reduce a evaluar
una fórmula cerrada (véase Sección~\ref{sec:extraccion_simplificado}).

El método simplificado es menos preciso a alta frecuencia y para sondas de mayor
apertura, pero es adecuado para exploraciones rápidas o cuando solo se dispone de
un líquido de referencia de composición conocida.

\newpage
\Encabezado

% -----------------------------------------------------------------------------
\subsection{Los Líquidos de Referencia: Selección y Criterios}
\label{sec:liquidos_referencia}
% -----------------------------------------------------------------------------

\subsubsection{Función de los líquidos de referencia}

Los líquidos de referencia cumplen una doble función:
\begin{itemize}
    \item \textbf{Calibración del analizador}: sus mediciones de $S_{11}$
          proveen puntos de la relación $\Gamma \leftrightarrow \varepsilon^*$
          con permitividades \emph{conocidas a priori} mediante los modelos
          termodinámicos de Debye.
    \item \textbf{Determinación de los parámetros de la sonda}: al conocer
          $\varepsilon^*$ se puede resolver el sistema de ecuaciones que da
          $C_0$ y $G_0$ sin necesidad de desmontarla ni medirla
          geométricamente.
\end{itemize}

\subsubsection{Por qué se necesitan dos líquidos en el método completo}

El modelo de admitancia~\eqref{eq:admitancia_modelo} tiene dos incógnitas
reales de geometría ($C_0$, $G_0$) más los tres errores del analizador,
totalizando cinco parámetros. Los estándares Open y Short aportan dos
ecuaciones, y cada líquido aporta una ecuación adicional (una por frecuencia,
pero los parámetros son constantes en frecuencia). Con un solo líquido se
tienen cuatro ecuaciones y cinco incógnitas: el sistema está indeterminado.
Con dos líquidos se completan las cinco ecuaciones necesarias.

Desde otro punto de vista: el primer líquido (Ref\,1 = IPA) fija principalmente
el coeficiente capacitivo~$C_0$, que domina la admitancia a baja frecuencia.
El segundo líquido (Ref\,2 = agua) —de permitividad considerablemente mayor—
fija principalmente el coeficiente de radiación~$G_0$, cuyo efecto resulta
apreciable a alta frecuencia y con materiales de alta permitividad.

\subsubsection{Por qué los líquidos deben ser diferentes entre sí}

Si se utilizaran dos líquidos de permitividad idéntica o muy similar,
sus mediciones de $S_{11}$ quedarían muy próximas en el plano complejo
y el sistema de ecuaciones resultaría mal condicionado numéricamente
(matriz casi singular). Esto se manifiesta como inestabilidad y ruido
amplificado en los valores de $C_0$ y $G_0$ extraídos.

Para un buen condicionamiento, los líquidos de referencia deben:
\begin{itemize}
    \item Tener partes reales $\varepsilon'$ significativamente distintas,
          idealmente separadas por al menos una década (por ejemplo,
          $\varepsilon'_{\text{IPA}} \approx 4$--$20$ frente a
          $\varepsilon'_{\text{agua}} \approx 40$--$80$ en la banda de 10~MHz a 3~GHz).
    \item Cubrir el rango de permitividades del material bajo prueba desde ambos lados
          siempre que sea posible.
    \item Tener permitividades tabuladas con alta confianza en el rango de temperatura
          y frecuencia de operación.
\end{itemize}

\subsubsection{Criterios de selección práctica}

Los líquidos de referencia ideales deben además satisfacer:
\begin{itemize}
    \item \textbf{Compatibilidad química}: no deben atacar el conector ni la
          brida de la sonda (acero inoxidable, PTFE, latón).
    \item \textbf{Bajo índice de evaporación}: la medición de $S_{11}$ tarda
          varios segundos; un líquido muy volátil puede cambiar su concentración
          durante la medición.
    \item \textbf{Homogeneidad}: deben ser líquidos puros, no mezclas cuya
          concentración sea difícil de controlar.
    \item \textbf{Datos de referencia disponibles}: se requieren coeficientes de Debye
          tabulados en función de temperatura, idealmente con trazabilidad a
          laboratorios nacionales de metrología (NPL, NIST, PTB).
\end{itemize}

En la Tabla~\ref{tab:liquidos_disponibles} se resumen las propiedades relevantes
de los cuatro líquidos disponibles en el sistema a 25~°C y 1~GHz (valores
aproximados; los valores exactos se calculan a partir de los coeficientes
tabulados).

\begin{table}[h!]
\centering
\caption{Propiedades dieléctricas aproximadas de los líquidos de referencia del sistema a 25~°C y 1~GHz.}
\label{tab:liquidos_disponibles}
\begin{tabular}{@{}lcccc@{}}
\toprule
\textbf{Líquido} & $\varepsilon'$ & $\varepsilon''$ & $\tan\delta$ & \textbf{Rol típico} \\
\midrule
Agua (H$_2$O)         & $\approx 78$  & $\approx 10$  & $\approx 0{,}13$ & Ref\,2 (alto $\varepsilon$) \\
Metanol               & $\approx 25$  & $\approx 17$  & $\approx 0{,}68$ & Ref\,2 alternativo \\
IPA (isopropanol)     & $\approx 15$  & $\approx 7$   & $\approx 0{,}47$ & Ref\,1 (bajo $\varepsilon$) \\
Etanol                & $\approx 22$  & $\approx 15$  & $\approx 0{,}68$ & Ref\,1 alternativo \\
\bottomrule
\end{tabular}
\end{table}

\noindent El par predeterminado del sistema es \textbf{IPA} (Ref\,1) + \textbf{agua} (Ref\,2),
combinación validada en la referencia de base~\cite{utnfrba2024} y ampliamente
utilizada en la literatura~\cite{stuchly1980, misra1990, gregory2012}.

\newpage
\Encabezado

% -----------------------------------------------------------------------------
\subsection{Algoritmo de Extracción de Permitividad}
\label{sec:extraccion}
% -----------------------------------------------------------------------------

\subsubsection{Método Completo: Polinomio de Grado 5}
\label{sec:extraccion_completo}

Una vez completada la calibración, los parámetros del sistema quedan condensados en
la \emph{constante de patrón} $G_n$ y en las admitancias normalizadas de los
líquidos de referencia~$y_1$, $y_2$.

\paragraph{Cálculo de $G_n$}

A partir de los cuatro $S_{11}$ de calibración y de las permitividades conocidas
$\varepsilon_1 = \varepsilon_{\text{Ref1}}$ y $\varepsilon_2 = \varepsilon_{\text{Ref2}}$,
se calculan los numeradores y denominadores:

\begin{align}
    \mathcal{A} &= (S_O - S_K)(S_2 - S_1), \nonumber \\
    \mathcal{B} &= (S_O - S_1)(S_K - S_2)\,\varepsilon_2, \nonumber \\
    \mathcal{C} &= (S_O - S_2)(S_1 - S_K)\,\varepsilon_1, \nonumber \\[4pt]
    \mathcal{D} &= (S_O - S_K)(S_2 - S_1), \nonumber \\
    \mathcal{E} &= (S_O - S_1)(S_K - S_2)\,\varepsilon_2^{5/2}, \nonumber \\
    \mathcal{F} &= (S_O - S_2)(S_1 - S_K)\,\varepsilon_1^{5/2}, \nonumber
\end{align}

\begin{equation}
    G_n = -\frac{\mathcal{A} + \mathcal{B} + \mathcal{C}}
                {\mathcal{D} + \mathcal{E} + \mathcal{F}},
    \label{eq:gn}
\end{equation}

\noindent donde $S_O$, $S_K$, $S_1$, $S_2$ son los coeficientes de reflexión
medidos para Open, Short, Ref\,1 y Ref\,2, respectivamente (todos complejos,
dependientes de la frecuencia).

Las admitancias normalizadas de los líquidos son:
\begin{align}
    y_1 &= \varepsilon_1 + G_n\,\varepsilon_1^{5/2}, \label{eq:y1} \\
    y_2 &= \varepsilon_2 + G_n\,\varepsilon_2^{5/2}. \label{eq:y2}
\end{align}

\paragraph{Término cruzado del material bajo prueba}

Dado el $S_{11}$ medido del material desconocido ($S_m$), se calcula un
\emph{término cruzado}~$T$ mediante la razón cruzada de los cuatro puntos
$\{S_m, S_K, S_1, S_2\}$:

\begin{equation}
    T = \frac{(S_m - S_1)(S_K - S_2)}{(S_m - S_K)(S_2 - S_1)}\,y_2
      + \frac{(S_m - S_2)(S_1 - S_K)}{(S_m - S_K)(S_2 - S_1)}\,y_1.
    \label{eq:third_term}
\end{equation}

\paragraph{Formulación del polinomio de grado 5}

Sustituyendo el modelo de admitancia~\eqref{eq:admitancia_modelo} normalizado
en la razón cruzada, y realizando el cambio de variable $x = \sqrt{\varepsilon_m^*}$,
se obtiene un polinomio de grado 5 en $x$:

\begin{equation}
    G_n\,x^5 + x^2 + T = 0,
    \label{eq:polinomio}
\end{equation}

que se resuelve numéricamente para cada punto en frecuencia, obteniendo cinco
raíces complejas. La permitividad se recupera como $\varepsilon_m^* = x^2$.

\paragraph{Selección de la raíz física}

De las cinco raíces del polinomio~\eqref{eq:polinomio}, solo aquellas con
\begin{equation}
    \operatorname{Re}(\varepsilon_m^*) > 0 \quad \text{y} \quad
    \operatorname{Im}(\varepsilon_m^*) \leq 0
    \label{eq:criterio_fisico}
\end{equation}
son físicamente admisibles (medio dieléctrico pasivo con la convención de
signo~\eqref{eq:permitividad_compleja}).

Cuando la condición~\eqref{eq:criterio_fisico} no resulta suficiente para
seleccionar una raíz unívoca, se aplica un algoritmo de \emph{seguimiento de
rama en frecuencia}: comenzando en la frecuencia más alta (donde los efectos
de radiación son máximos y la raíz física es más distinguible), se sigue la
rama de raíces de mayor continuidad hacia frecuencias decrecientes, apoyándose
en una extrapolación polinómica de orden~2 con ventana de 5 puntos.

Como referencia inicial se utiliza el resultado del método simplificado
(Sección~\ref{sec:extraccion_simplificado}), que proporciona una solución
aproximada sin ambigüedad de raíces.

\paragraph{Verificación de consistencia}

Al finalizar la extracción, el sistema compara el resultado del método completo
con el del simplificado en la mitad inferior de la banda. Si la divergencia
mediana supera el 25\,\%, se emite una advertencia al operador indicando posible
mal acondicionamiento de la calibración o elección inadecuada de líquidos.

\subsubsection{Método Simplificado: Transformada Bilineal}
\label{sec:extraccion_simplificado}

Con $G_n = 0$ (sin término de radiación), el modelo de admitancia se reduce a
$Y = j\omega C_0\,\varepsilon^*$, y la relación entre $S_{11}$ y $\varepsilon^*$ es
una transformada de Möbius exacta. Cuatro puntos en el plano~$S_{11}$ y sus
correspondientes en el plano~$\varepsilon$ determinan la razón cruzada.

Identificando Open ($\varepsilon_{\text{aire}} \approx 1$), Short
($\varepsilon_{\text{cortocircuito}} \to \infty$, contribución límite) y
Ref\,1 ($\varepsilon_{\text{Ref}}$), la permitividad del material desconocido
se calcula directamente como:

\begin{equation}
    \varepsilon_m^*
    = -\,c_{\text{Ref}}\,\varepsilon_{\text{Ref}}^*
      -\,c_{\text{aire}}\,\varepsilon_{\text{aire}},
    \label{eq:formula_simplificada}
\end{equation}

\noindent con coeficientes bariocéntricos:

\begin{align}
    c_{\text{Ref}}  &= \frac{(S_m - S_O)(S_K - S_{\text{Ref}})}{(S_m - S_K)(S_{\text{Ref}} - S_O)},
    \label{eq:c_ref} \\[4pt]
    c_{\text{aire}} &= \frac{(S_m - S_{\text{Ref}})(S_O - S_K)}{(S_m - S_K)(S_{\text{Ref}} - S_O)}.
    \label{eq:c_aire}
\end{align}

La identidad $c_{\text{Ref}} + c_{\text{aire}} = -1$ es una consecuencia
directa de la razón cruzada (identidad de Ptolemy), y sirve como verificación
numérica de la coherencia del cálculo.

Nótese que~\eqref{eq:formula_simplificada} es algebraicamente cerrada: no
requiere resolver polinomios ni realizar seguimiento de ramas, lo que la hace
computacionalmente trivial y robusta frente a señales ruidosas.

\newpage
\Encabezado

% -----------------------------------------------------------------------------
\subsection{Flujo Completo del Algoritmo}
\label{sec:flujo}
% -----------------------------------------------------------------------------

La Figura~\ref{fig:flujo_algoritmo} (pendiente de inclusión) resume el flujo
de datos del sistema. En forma textual, las etapas son:

\begin{enumerate}
    \item \textbf{Configuración}: selección de sonda ($C_0$ tabulated por modelo),
          temperatura de medición, líquidos de referencia y método
          (completo o simplificado).

    \item \textbf{Fase de calibración}:
    \begin{itemize}
        \item Medir o cargar los valores presetados de $S_{11}$ para Open, Short,
              Ref\,1 (y Ref\,2 en método completo).
        \item Evaluar $\varepsilon_{\text{Ref}}^*(f, T)$ mediante los modelos de
              Debye de la Sección~\ref{sec:debye} interpolando en temperatura.
        \item Calcular $G_n(f)$, $y_1(f)$, $y_2(f)$
              según~\eqref{eq:gn}--\eqref{eq:y2} (solo método completo).
    \end{itemize}

    \item \textbf{Medición del material bajo prueba}: barrido en frecuencia con
          la NanoVNA para obtener $S_{11,m}(f)$.

    \item \textbf{Extracción de $\varepsilon^*(f)$}:
    \begin{itemize}
        \item \emph{Método completo}: calcular $T(f)$ mediante~\eqref{eq:third_term},
              resolver el polinomio~\eqref{eq:polinomio} en cada frecuencia,
              seleccionar la raíz física con seguimiento de rama.
        \item \emph{Método simplificado}: evaluar directamente~\eqref{eq:formula_simplificada}.
    \end{itemize}

    \item \textbf{Salida}: curvas $\varepsilon'(f)$ y $\varepsilon''(f)$
          presentadas en la pantalla de resultados. El resultado puede exportarse
          como archivo Touchstone~`.s1p` para análisis externo.
\end{enumerate}

\newpage
\Encabezado

% -----------------------------------------------------------------------------
\subsection{Resumen Comparativo de los Métodos}
\label{sec:resumen}
% -----------------------------------------------------------------------------

\begin{table}[h!]
\centering
\caption{Comparación entre el método completo y el método simplificado.}
\label{tab:comparacion_metodos}
\begin{tabular}{@{}lll@{}}
\toprule
\textbf{Característica}         & \textbf{Método completo}              & \textbf{Método simplificado} \\
\midrule
Estándares de calibración       & 4 (Open, Short, Ref1, Ref2)           & 3 (Open, Short, Ref) \\
Líquidos de referencia          & 2 (IPA + Agua)                        & 1 (Agua) \\
Modelo de admitancia            & $j\omega C_0\varepsilon + G_0\varepsilon^{5/2}$ & $j\omega C_0\varepsilon$ \\
Extracción                      & Polinomio grado 5 + seguimiento       & Fórmula cerrada bilineal \\
Precisión a alta frecuencia     & Alta                                  & Moderada \\
Precisión para sonda de 21\,mm  & Alta                                  & Reducida \\
Complejidad computacional       & Media                                 & Baja \\
Velocidad de ensayo             & Normal                                & Rápida (exploración) \\
\bottomrule
\end{tabular}
\end{table}

% -----------------------------------------------------------------------------
%  REFERENCIAS
% -----------------------------------------------------------------------------
% (Agregar al archivo .bib del informe o usar directamente \bibitem)
%
% \begin{thebibliography}{9}
%
% \bibitem{debye1929}
% P.~Debye,
% \emph{Polar Molecules}.
% Chemical Catalog Company, New York, 1929.
%
% \bibitem{stuchly1980}
% S.~S.~Stuchly y M.~A.~Stuchly,
% ``Coaxial line reflection methods for measuring dielectric properties of biological
%   substances at radio and microwave frequencies — A review,''
% \emph{IEEE Trans. Instrum. Meas.}, vol.~29, no.~3, pp.~176--183, Sep. 1980.
%
% \bibitem{misra1990}
% D.~Misra,
% ``A quasi-static analysis of open-ended coaxial lines,''
% \emph{IEEE Trans. Microw. Theory Techn.}, vol.~35, no.~10, pp.~925--928, Oct. 1987.
%
% \bibitem{gregory2012}
% A.~P.~Gregory y R.~N.~Clarke,
% \emph{A review of RF and microwave techniques for dielectric measurements on polar liquids},
% NPL Report MAT 23. National Physical Laboratory, Teddington, UK, 2012.
%
% \bibitem{utnfrba2024}
% [COMPLETAR: autores],
% ``Mediciones de permitividad mediante una sonda coaxial,''
% UTN-FRBA, 2024. (Informe interno de referencia del proyecto.)
%
% \end{thebibliography}

%    \end{document}
% =============================================================================

\section{Marco Teórico: Caracterización Dieléctrica mediante Sonda Coaxial Abierta}

El presente capítulo desarrolla los fundamentos físicos, matemáticos y metrológicos
del método de sonda coaxial abierta empleado por el sistema NanoVNA-UTN-Toolkit
para la determinación de la permitividad eléctrica compleja de materiales líquidos.
Se presentan los dos algoritmos implementados —el método completo y el simplificado—,
los modelos de Debye utilizados para los líquidos de referencia, y los criterios de
selección de dichos líquidos.

\newpage
\Encabezado

% -----------------------------------------------------------------------------
\subsection{La Permitividad Eléctrica Compleja}
% -----------------------------------------------------------------------------

La permitividad eléctrica $\varepsilon$ es la magnitud que describe la respuesta de
un material a la aplicación de un campo eléctrico externo. En el dominio frecuencial,
para un material dieléctrico con pérdidas, se expresa como un número complejo:

\begin{equation}
    \varepsilon^*(\omega) = \varepsilon'(\omega) - j\,\varepsilon''(\omega),
    \label{eq:permitividad_compleja}
\end{equation}

\noindent donde:
\begin{itemize}
    \item $\varepsilon'$ es la parte real, asociada a la energía almacenada en el campo
          eléctrico (capacitiva). Cuantifica cuánto puede polarizarse el material.
    \item $\varepsilon''$ es la parte imaginaria, siempre no negativa para un medio
          pasivo, asociada a las pérdidas dieléctricas (conversión de energía
          electromagnética en calor).
    \item $\omega = 2\pi f$ es la frecuencia angular.
\end{itemize}

La permitividad relativa $\varepsilon_r$ se define normalizando respecto a la
permitividad del vacío $\varepsilon_0 = 8{,}854 \times 10^{-12}$~F/m:

\begin{equation}
    \varepsilon_r^* = \frac{\varepsilon^*}{\varepsilon_0}
                    = \varepsilon_r' - j\,\varepsilon_r''.
\end{equation}

Una magnitud derivada de uso frecuente es el \emph{factor de disipación} o
\emph{tangente de pérdidas}:

\begin{equation}
    \tan\delta = \frac{\varepsilon_r''}{\varepsilon_r'}.
\end{equation}

A lo largo de este documento se omite el subíndice $r$ salvo aclaración explícita;
todas las magnitudes son relativas.

\newpage
\Encabezado

% -----------------------------------------------------------------------------
\subsection{Modelo de Debye para Líquidos Polares}
\label{sec:debye}
% -----------------------------------------------------------------------------

El modelo de Debye~\cite{debye1929} describe la dependencia en frecuencia de la
permitividad de un dieléctrico polar cuyas moléculas poseen un único tiempo de
relajación $\tau$. Físicamente, $\tau$ cuantifica el tiempo característico de
reorientación dipolar ante un campo oscilante.

\subsubsection{Modelo de Debye simple}

\begin{equation}
    \varepsilon^*(f) = \varepsilon_\infty
                     + \frac{\varepsilon_s - \varepsilon_\infty}{1 + j\,f / f_r},
    \label{eq:debye_simple}
\end{equation}

\noindent donde:
\begin{itemize}
    \item $\varepsilon_s$ es la permitividad estática (límite $f \to 0$).
    \item $\varepsilon_\infty$ es la permitividad a alta frecuencia (límite $f \to \infty$),
          debida a polarizaciones electrónicas e iónicas.
    \item $f_r = 1/(2\pi\tau)$ es la frecuencia de relajación.
\end{itemize}

Este modelo se aplica al agua y al metanol en el presente sistema.
Los coeficientes $\{\varepsilon_s, \varepsilon_\infty, \tau\}$ dependen de la
temperatura y se tabulan a partir de los datos del NPL Report MAT~23
\cite{gregory2012} mediante ajuste polinomial cuadrático en temperatura.

\subsubsection{Modelo de Debye doble}

Algunos líquidos exhiben dos mecanismos de relajación dominantes a frecuencias
distintas (por ejemplo, el isopropanol / IPA). Se modela con dos términos de Debye
superpuestos:

\begin{equation}
    \varepsilon^*(f) = \varepsilon_\infty
                     + \frac{\varepsilon_s - \varepsilon_h}{1 + j\,f / f_{r1}}
                     + \frac{\varepsilon_h - \varepsilon_\infty}{1 + j\,f / f_{r2}},
    \label{eq:debye_doble}
\end{equation}

\noindent donde $\varepsilon_h$ es una permitividad intermedia y $f_{r1} < f_{r2}$
son las dos frecuencias de relajación.

\subsubsection{Modelo de Debye-Gamma}

Para el etanol, el NPL Report MAT~23 propone un modelo de Debye simple con
corrección lineal en frecuencia (término~$\gamma$) que aproxima una segunda
relajación en la parte superior de la banda:

\begin{equation}
    \varepsilon^*(f) = \varepsilon_h
                     + \frac{\varepsilon_s - \varepsilon_h}{1 + j\,f / f_r}
                     - j\,\gamma \left(\frac{f}{\text{GHz}}\right),
    \label{eq:debye_gamma}
\end{equation}

siendo $\gamma$ un coeficiente dimensional en GHz$^{-1}$.

En la Tabla~\ref{tab:modelos_debye} se resume el modelo utilizado para cada
líquido disponible en el sistema.

\begin{table}[h!]
\centering
\caption{Modelos de Debye por líquido de referencia implementados en el sistema.}
\label{tab:modelos_debye}
\begin{tabular}{@{}llcc@{}}
\toprule
\textbf{Líquido} & \textbf{Modelo} & \textbf{Rango T (°C)} & \textbf{Fuente} \\
\midrule
Agua (H$_2$O) & Debye simple       & 0 -- 60  & NPL MAT 23 \cite{gregory2012} \\
IPA (isopropanol) & Debye doble    & 10 -- 30 & NPL MAT 23 \cite{gregory2012} \\
Etanol         & Debye-Gamma        & 10 -- 50 & NPL MAT 23 \cite{gregory2012} \\
Metanol        & Debye simple       & 10 -- 50 & NPL MAT 23 \cite{gregory2012} \\
\bottomrule
\end{tabular}
\end{table}

\newpage
\Encabezado

% -----------------------------------------------------------------------------
\subsection{Método de la Sonda Coaxial Abierta}
\label{sec:sonda_coaxial}
% -----------------------------------------------------------------------------

\subsubsection{Principio físico}

La técnica de sonda coaxial abierta es un método de reflexión no resonante
y no destructivo para la medición de permitividad~\cite{stuchly1980, misra1990}.
Consiste en apoyar o sumergir el extremo abierto de una línea coaxial sobre
el material bajo prueba (MBP) y medir el coeficiente de reflexión~$\Gamma$
(parámetro~$S_{11}$) en un rango de frecuencias.

Cuando el extremo de la línea coaxial se abre a un semiespacio relleno con
el material, la terminación coaxial abierta actúa como una antena de apertura
en campo cercano. La admitancia en la terminación depende fuertemente de las
propiedades dieléctricas del material en contacto:

\begin{equation}
    Y(\varepsilon^*, \omega)
    = j\omega C_0\,\varepsilon^*
    + G_0\,(\varepsilon^*)^{5/2},
    \label{eq:admitancia_modelo}
\end{equation}

\noindent donde:
\begin{itemize}
    \item $C_0$ [F] es la capacitancia de la sonda en el vacío, determinada por la
          geometría del coaxial. Domina el comportamiento a bajas frecuencias.
    \item $G_0$ [S] es el coeficiente de radiación. El término $G_0\,\varepsilon^{5/2}$
          representa las pérdidas por radiación en el campo lejano; domina a altas
          frecuencias y para sondas de mayor apertura (sonda de 21~mm respecto
          a la de 3~mm).
\end{itemize}

La relación entre la admitancia y el coeficiente de reflexión medido es la
transformada de Möbius estándar de la teoría de redes:

\begin{equation}
    \Gamma = \frac{Z_0\,Y - 1}{Z_0\,Y + 1}
    \quad \Longleftrightarrow \quad
    Y = \frac{1}{Z_0}\cdot\frac{1 - \Gamma}{1 + \Gamma},
\end{equation}

\noindent con $Z_0 = 50~\Omega$.

\subsubsection{Historia y contexto del método}

El método de sonda coaxial abierta fue formalizado matemáticamente en la
década de 1980 por Stuchly et al.~\cite{stuchly1980} y Misra~\cite{misra1990},
quienes propusieron el modelo de admitancia de~\eqref{eq:admitancia_modelo}
y demostraron que el exponente $5/2$ emerge del desarrollo en serie del campo
radiado por una apertura circular en un plano conductor infinito.

La calibración con líquidos de referencia, en lugar de estándares mecánicos de
reflexión, fue introducida para eliminar la dependencia del resultado respecto
a la calidad del contacto sonda-muestra y respecto a los errores sistemáticos
del analizador de redes. La versión de dos líquidos que se documenta aquí es
consistente con la referencia UTN-FRBA 2024~\cite{utnfrba2024} en la que se basa
el presente trabajo.

\newpage
\Encabezado

% -----------------------------------------------------------------------------
\subsection{Calibración del Sistema}
\label{sec:calibracion}
% -----------------------------------------------------------------------------

El sistema presenta dos métodos de calibración con distinto número de estándares,
distintos supuestos y distinta complejidad algorítmica.

\subsubsection{Método Completo (4 estándares)}
\label{sec:metodo_completo}

Este método es denominado en la literatura como calibración de 4 estándares
del tipo SOL + 2 líquidos (Open -- Short -- Liquid\,1 -- Liquid\,2).
Los cuatro estándares son:

\begin{enumerate}
    \item \textbf{Circuito abierto (Open)}: la sonda apunta al aire.
          $\varepsilon_{\text{open}} = \varepsilon_{\text{aire}} \approx 1$.
    \item \textbf{Cortocircuito (Short)}: la apertura de la sonda se cubre con
          una lámina metálica conductora.
    \item \textbf{Líquido de referencia~1 (Ref\,1)}: por defecto IPA (isopropanol),
          de permitividad baja-moderada.
    \item \textbf{Líquido de referencia~2 (Ref\,2)}: por defecto agua destilada,
          de permitividad alta.
\end{enumerate}

Las cuatro mediciones $S_{11}$ de calibración y las permitividades conocidas
$\varepsilon_{\text{Ref1}}$, $\varepsilon_{\text{Ref2}}$ (calculadas mediante
los modelos de Debye de la Sección~\ref{sec:debye} a la temperatura de medición)
permiten determinar los parámetros del sistema $\{C_0, G_0\}$ y los tres
coeficientes de error de la caja de errores del analizador (directividad,
adaptación de fuente, seguimiento de reflexión).

La calibración completa absorbe todos los errores sistemáticos del analizador y
del cable de la sonda simultáneamente con los parámetros de la sonda misma.
No se requiere una calibración OSM previa del analizador.

\subsubsection{Método Simplificado (3 estándares)}
\label{sec:metodo_simplificado}

El método simplificado (denominado también \emph{incompleto} en la literatura)
utiliza únicamente tres estándares:

\begin{enumerate}
    \item \textbf{Open}: ídem al método completo.
    \item \textbf{Short}: ídem al método completo.
    \item \textbf{Líquido de referencia único (Ref\,1)}: por defecto agua destilada.
\end{enumerate}

La simplificación consiste en despreciar el término de radiación ($G_0 = 0$),
reduciendo el modelo de admitancia a uno puramente capacitivo:

\begin{equation}
    Y \approx j\omega C_0\,\varepsilon^*.
\end{equation}

Con esta hipótesis, la relación entre $S_{11}$ y $\varepsilon^*$ se vuelve una
transformada de Möbius (bilineal) exacta. Tres puntos conocidos en el plano~$S_{11}$
(aire, cortocircuito, líquido de referencia) determinan unívocamente la transformada,
y la extracción de $\varepsilon^*$ del material desconocido se reduce a evaluar
una fórmula cerrada (véase Sección~\ref{sec:extraccion_simplificado}).

El método simplificado es menos preciso a alta frecuencia y para sondas de mayor
apertura, pero es adecuado para exploraciones rápidas o cuando solo se dispone de
un líquido de referencia de composición conocida.

\newpage
\Encabezado

% -----------------------------------------------------------------------------
\subsection{Los Líquidos de Referencia: Selección y Criterios}
\label{sec:liquidos_referencia}
% -----------------------------------------------------------------------------

\subsubsection{Función de los líquidos de referencia}

Los líquidos de referencia cumplen una doble función:
\begin{itemize}
    \item \textbf{Calibración del analizador}: sus mediciones de $S_{11}$
          proveen puntos de la relación $\Gamma \leftrightarrow \varepsilon^*$
          con permitividades \emph{conocidas a priori} mediante los modelos
          termodinámicos de Debye.
    \item \textbf{Determinación de los parámetros de la sonda}: al conocer
          $\varepsilon^*$ se puede resolver el sistema de ecuaciones que da
          $C_0$ y $G_0$ sin necesidad de desmontarla ni medirla
          geométricamente.
\end{itemize}

\subsubsection{Por qué se necesitan dos líquidos en el método completo}

El modelo de admitancia~\eqref{eq:admitancia_modelo} tiene dos incógnitas
reales de geometría ($C_0$, $G_0$) más los tres errores del analizador,
totalizando cinco parámetros. Los estándares Open y Short aportan dos
ecuaciones, y cada líquido aporta una ecuación adicional (una por frecuencia,
pero los parámetros son constantes en frecuencia). Con un solo líquido se
tienen cuatro ecuaciones y cinco incógnitas: el sistema está indeterminado.
Con dos líquidos se completan las cinco ecuaciones necesarias.

Desde otro punto de vista: el primer líquido (Ref\,1 = IPA) fija principalmente
el coeficiente capacitivo~$C_0$, que domina la admitancia a baja frecuencia.
El segundo líquido (Ref\,2 = agua) —de permitividad considerablemente mayor—
fija principalmente el coeficiente de radiación~$G_0$, cuyo efecto resulta
apreciable a alta frecuencia y con materiales de alta permitividad.

\subsubsection{Por qué los líquidos deben ser diferentes entre sí}

Si se utilizaran dos líquidos de permitividad idéntica o muy similar,
sus mediciones de $S_{11}$ quedarían muy próximas en el plano complejo
y el sistema de ecuaciones resultaría mal condicionado numéricamente
(matriz casi singular). Esto se manifiesta como inestabilidad y ruido
amplificado en los valores de $C_0$ y $G_0$ extraídos.

Para un buen condicionamiento, los líquidos de referencia deben:
\begin{itemize}
    \item Tener partes reales $\varepsilon'$ significativamente distintas,
          idealmente separadas por al menos una década (por ejemplo,
          $\varepsilon'_{\text{IPA}} \approx 4$--$20$ frente a
          $\varepsilon'_{\text{agua}} \approx 40$--$80$ en la banda de 10~MHz a 3~GHz).
    \item Cubrir el rango de permitividades del material bajo prueba desde ambos lados
          siempre que sea posible.
    \item Tener permitividades tabuladas con alta confianza en el rango de temperatura
          y frecuencia de operación.
\end{itemize}

\subsubsection{Criterios de selección práctica}

Los líquidos de referencia ideales deben además satisfacer:
\begin{itemize}
    \item \textbf{Compatibilidad química}: no deben atacar el conector ni la
          brida de la sonda (acero inoxidable, PTFE, latón).
    \item \textbf{Bajo índice de evaporación}: la medición de $S_{11}$ tarda
          varios segundos; un líquido muy volátil puede cambiar su concentración
          durante la medición.
    \item \textbf{Homogeneidad}: deben ser líquidos puros, no mezclas cuya
          concentración sea difícil de controlar.
    \item \textbf{Datos de referencia disponibles}: se requieren coeficientes de Debye
          tabulados en función de temperatura, idealmente con trazabilidad a
          laboratorios nacionales de metrología (NPL, NIST, PTB).
\end{itemize}

En la Tabla~\ref{tab:liquidos_disponibles} se resumen las propiedades relevantes
de los cuatro líquidos disponibles en el sistema a 25~°C y 1~GHz (valores
aproximados; los valores exactos se calculan a partir de los coeficientes
tabulados).

\begin{table}[h!]
\centering
\caption{Propiedades dieléctricas aproximadas de los líquidos de referencia del sistema a 25~°C y 1~GHz.}
\label{tab:liquidos_disponibles}
\begin{tabular}{@{}lcccc@{}}
\toprule
\textbf{Líquido} & $\varepsilon'$ & $\varepsilon''$ & $\tan\delta$ & \textbf{Rol típico} \\
\midrule
Agua (H$_2$O)         & $\approx 78$  & $\approx 10$  & $\approx 0{,}13$ & Ref\,2 (alto $\varepsilon$) \\
Metanol               & $\approx 25$  & $\approx 17$  & $\approx 0{,}68$ & Ref\,2 alternativo \\
IPA (isopropanol)     & $\approx 15$  & $\approx 7$   & $\approx 0{,}47$ & Ref\,1 (bajo $\varepsilon$) \\
Etanol                & $\approx 22$  & $\approx 15$  & $\approx 0{,}68$ & Ref\,1 alternativo \\
\bottomrule
\end{tabular}
\end{table}

\noindent El par predeterminado del sistema es \textbf{IPA} (Ref\,1) + \textbf{agua} (Ref\,2),
combinación validada en la referencia de base~\cite{utnfrba2024} y ampliamente
utilizada en la literatura~\cite{stuchly1980, misra1990, gregory2012}.

\newpage
\Encabezado

% -----------------------------------------------------------------------------
\subsection{Algoritmo de Extracción de Permitividad}
\label{sec:extraccion}
% -----------------------------------------------------------------------------

\subsubsection{Método Completo: Polinomio de Grado 5}
\label{sec:extraccion_completo}

Una vez completada la calibración, los parámetros del sistema quedan condensados en
la \emph{constante de patrón} $G_n$ y en las admitancias normalizadas de los
líquidos de referencia~$y_1$, $y_2$.

\paragraph{Cálculo de $G_n$}

A partir de los cuatro $S_{11}$ de calibración y de las permitividades conocidas
$\varepsilon_1 = \varepsilon_{\text{Ref1}}$ y $\varepsilon_2 = \varepsilon_{\text{Ref2}}$,
se calculan los numeradores y denominadores:

\begin{align}
    \mathcal{A} &= (S_O - S_K)(S_2 - S_1), \nonumber \\
    \mathcal{B} &= (S_O - S_1)(S_K - S_2)\,\varepsilon_2, \nonumber \\
    \mathcal{C} &= (S_O - S_2)(S_1 - S_K)\,\varepsilon_1, \nonumber \\[4pt]
    \mathcal{D} &= (S_O - S_K)(S_2 - S_1), \nonumber \\
    \mathcal{E} &= (S_O - S_1)(S_K - S_2)\,\varepsilon_2^{5/2}, \nonumber \\
    \mathcal{F} &= (S_O - S_2)(S_1 - S_K)\,\varepsilon_1^{5/2}, \nonumber
\end{align}

\begin{equation}
    G_n = -\frac{\mathcal{A} + \mathcal{B} + \mathcal{C}}
                {\mathcal{D} + \mathcal{E} + \mathcal{F}},
    \label{eq:gn}
\end{equation}

\noindent donde $S_O$, $S_K$, $S_1$, $S_2$ son los coeficientes de reflexión
medidos para Open, Short, Ref\,1 y Ref\,2, respectivamente (todos complejos,
dependientes de la frecuencia).

Las admitancias normalizadas de los líquidos son:
\begin{align}
    y_1 &= \varepsilon_1 + G_n\,\varepsilon_1^{5/2}, \label{eq:y1} \\
    y_2 &= \varepsilon_2 + G_n\,\varepsilon_2^{5/2}. \label{eq:y2}
\end{align}

\paragraph{Término cruzado del material bajo prueba}

Dado el $S_{11}$ medido del material desconocido ($S_m$), se calcula un
\emph{término cruzado}~$T$ mediante la razón cruzada de los cuatro puntos
$\{S_m, S_K, S_1, S_2\}$:

\begin{equation}
    T = \frac{(S_m - S_1)(S_K - S_2)}{(S_m - S_K)(S_2 - S_1)}\,y_2
      + \frac{(S_m - S_2)(S_1 - S_K)}{(S_m - S_K)(S_2 - S_1)}\,y_1.
    \label{eq:third_term}
\end{equation}

\paragraph{Formulación del polinomio de grado 5}

Sustituyendo el modelo de admitancia~\eqref{eq:admitancia_modelo} normalizado
en la razón cruzada, y realizando el cambio de variable $x = \sqrt{\varepsilon_m^*}$,
se obtiene un polinomio de grado 5 en $x$:

\begin{equation}
    G_n\,x^5 + x^2 + T = 0,
    \label{eq:polinomio}
\end{equation}

que se resuelve numéricamente para cada punto en frecuencia, obteniendo cinco
raíces complejas. La permitividad se recupera como $\varepsilon_m^* = x^2$.

\paragraph{Selección de la raíz física}

De las cinco raíces del polinomio~\eqref{eq:polinomio}, solo aquellas con
\begin{equation}
    \operatorname{Re}(\varepsilon_m^*) > 0 \quad \text{y} \quad
    \operatorname{Im}(\varepsilon_m^*) \leq 0
    \label{eq:criterio_fisico}
\end{equation}
son físicamente admisibles (medio dieléctrico pasivo con la convención de
signo~\eqref{eq:permitividad_compleja}).

Cuando la condición~\eqref{eq:criterio_fisico} no resulta suficiente para
seleccionar una raíz unívoca, se aplica un algoritmo de \emph{seguimiento de
rama en frecuencia}: comenzando en la frecuencia más alta (donde los efectos
de radiación son máximos y la raíz física es más distinguible), se sigue la
rama de raíces de mayor continuidad hacia frecuencias decrecientes, apoyándose
en una extrapolación polinómica de orden~2 con ventana de 5 puntos.

Como referencia inicial se utiliza el resultado del método simplificado
(Sección~\ref{sec:extraccion_simplificado}), que proporciona una solución
aproximada sin ambigüedad de raíces.

\paragraph{Verificación de consistencia}

Al finalizar la extracción, el sistema compara el resultado del método completo
con el del simplificado en la mitad inferior de la banda. Si la divergencia
mediana supera el 25\,\%, se emite una advertencia al operador indicando posible
mal acondicionamiento de la calibración o elección inadecuada de líquidos.

\subsubsection{Método Simplificado: Transformada Bilineal}
\label{sec:extraccion_simplificado}

Con $G_n = 0$ (sin término de radiación), el modelo de admitancia se reduce a
$Y = j\omega C_0\,\varepsilon^*$, y la relación entre $S_{11}$ y $\varepsilon^*$ es
una transformada de Möbius exacta. Cuatro puntos en el plano~$S_{11}$ y sus
correspondientes en el plano~$\varepsilon$ determinan la razón cruzada.

Identificando Open ($\varepsilon_{\text{aire}} \approx 1$), Short
($\varepsilon_{\text{cortocircuito}} \to \infty$, contribución límite) y
Ref\,1 ($\varepsilon_{\text{Ref}}$), la permitividad del material desconocido
se calcula directamente como:

\begin{equation}
    \varepsilon_m^*
    = -\,c_{\text{Ref}}\,\varepsilon_{\text{Ref}}^*
      -\,c_{\text{aire}}\,\varepsilon_{\text{aire}},
    \label{eq:formula_simplificada}
\end{equation}

\noindent con coeficientes bariocéntricos:

\begin{align}
    c_{\text{Ref}}  &= \frac{(S_m - S_O)(S_K - S_{\text{Ref}})}{(S_m - S_K)(S_{\text{Ref}} - S_O)},
    \label{eq:c_ref} \\[4pt]
    c_{\text{aire}} &= \frac{(S_m - S_{\text{Ref}})(S_O - S_K)}{(S_m - S_K)(S_{\text{Ref}} - S_O)}.
    \label{eq:c_aire}
\end{align}

La identidad $c_{\text{Ref}} + c_{\text{aire}} = -1$ es una consecuencia
directa de la razón cruzada (identidad de Ptolemy), y sirve como verificación
numérica de la coherencia del cálculo.

Nótese que~\eqref{eq:formula_simplificada} es algebraicamente cerrada: no
requiere resolver polinomios ni realizar seguimiento de ramas, lo que la hace
computacionalmente trivial y robusta frente a señales ruidosas.

\newpage
\Encabezado

% -----------------------------------------------------------------------------
\subsection{Flujo Completo del Algoritmo}
\label{sec:flujo}
% -----------------------------------------------------------------------------

La Figura~\ref{fig:flujo_algoritmo} (pendiente de inclusión) resume el flujo
de datos del sistema. En forma textual, las etapas son:

\begin{enumerate}
    \item \textbf{Configuración}: selección de sonda ($C_0$ tabulated por modelo),
          temperatura de medición, líquidos de referencia y método
          (completo o simplificado).

    \item \textbf{Fase de calibración}:
    \begin{itemize}
        \item Medir o cargar los valores presetados de $S_{11}$ para Open, Short,
              Ref\,1 (y Ref\,2 en método completo).
        \item Evaluar $\varepsilon_{\text{Ref}}^*(f, T)$ mediante los modelos de
              Debye de la Sección~\ref{sec:debye} interpolando en temperatura.
        \item Calcular $G_n(f)$, $y_1(f)$, $y_2(f)$
              según~\eqref{eq:gn}--\eqref{eq:y2} (solo método completo).
    \end{itemize}

    \item \textbf{Medición del material bajo prueba}: barrido en frecuencia con
          la NanoVNA para obtener $S_{11,m}(f)$.

    \item \textbf{Extracción de $\varepsilon^*(f)$}:
    \begin{itemize}
        \item \emph{Método completo}: calcular $T(f)$ mediante~\eqref{eq:third_term},
              resolver el polinomio~\eqref{eq:polinomio} en cada frecuencia,
              seleccionar la raíz física con seguimiento de rama.
        \item \emph{Método simplificado}: evaluar directamente~\eqref{eq:formula_simplificada}.
    \end{itemize}

    \item \textbf{Salida}: curvas $\varepsilon'(f)$ y $\varepsilon''(f)$
          presentadas en la pantalla de resultados. El resultado puede exportarse
          como archivo Touchstone~`.s1p` para análisis externo.
\end{enumerate}

\newpage
\Encabezado

% -----------------------------------------------------------------------------
\subsection{Resumen Comparativo de los Métodos}
\label{sec:resumen}
% -----------------------------------------------------------------------------

\begin{table}[h!]
\centering
\caption{Comparación entre el método completo y el método simplificado.}
\label{tab:comparacion_metodos}
\begin{tabular}{@{}lll@{}}
\toprule
\textbf{Característica}         & \textbf{Método completo}              & \textbf{Método simplificado} \\
\midrule
Estándares de calibración       & 4 (Open, Short, Ref1, Ref2)           & 3 (Open, Short, Ref) \\
Líquidos de referencia          & 2 (IPA + Agua)                        & 1 (Agua) \\
Modelo de admitancia            & $j\omega C_0\varepsilon + G_0\varepsilon^{5/2}$ & $j\omega C_0\varepsilon$ \\
Extracción                      & Polinomio grado 5 + seguimiento       & Fórmula cerrada bilineal \\
Precisión a alta frecuencia     & Alta                                  & Moderada \\
Precisión para sonda de 21\,mm  & Alta                                  & Reducida \\
Complejidad computacional       & Media                                 & Baja \\
Velocidad de ensayo             & Normal                                & Rápida (exploración) \\
\bottomrule
\end{tabular}
\end{table}

% -----------------------------------------------------------------------------
%  REFERENCIAS
% -----------------------------------------------------------------------------
% (Agregar al archivo .bib del informe o usar directamente \bibitem)
%
% \begin{thebibliography}{9}
%
% \bibitem{debye1929}
% P.~Debye,
% \emph{Polar Molecules}.
% Chemical Catalog Company, New York, 1929.
%
% \bibitem{stuchly1980}
% S.~S.~Stuchly y M.~A.~Stuchly,
% ``Coaxial line reflection methods for measuring dielectric properties of biological
%   substances at radio and microwave frequencies — A review,''
% \emph{IEEE Trans. Instrum. Meas.}, vol.~29, no.~3, pp.~176--183, Sep. 1980.
%
% \bibitem{misra1990}
% D.~Misra,
% ``A quasi-static analysis of open-ended coaxial lines,''
% \emph{IEEE Trans. Microw. Theory Techn.}, vol.~35, no.~10, pp.~925--928, Oct. 1987.
%
% \bibitem{gregory2012}
% A.~P.~Gregory y R.~N.~Clarke,
% \emph{A review of RF and microwave techniques for dielectric measurements on polar liquids},
% NPL Report MAT 23. National Physical Laboratory, Teddington, UK, 2012.
%
% \bibitem{utnfrba2024}
% [COMPLETAR: autores],
% ``Mediciones de permitividad mediante una sonda coaxial,''
% UTN-FRBA, 2024. (Informe interno de referencia del proyecto.)
%
% \end{thebibliography}

%    \end{document}
% =============================================================================

\section{Marco Teórico: Caracterización Dieléctrica mediante Sonda Coaxial Abierta}

El presente capítulo desarrolla los fundamentos físicos, matemáticos y metrológicos
del método de sonda coaxial abierta empleado por el sistema NanoVNA-UTN-Toolkit
para la determinación de la permitividad eléctrica compleja de materiales líquidos.
Se presentan los dos algoritmos implementados —el método completo y el simplificado—,
los modelos de Debye utilizados para los líquidos de referencia, y los criterios de
selección de dichos líquidos.

\newpage
\Encabezado

% -----------------------------------------------------------------------------
\subsection{La Permitividad Eléctrica Compleja}
% -----------------------------------------------------------------------------

La permitividad eléctrica $\varepsilon$ es la magnitud que describe la respuesta de
un material a la aplicación de un campo eléctrico externo. En el dominio frecuencial,
para un material dieléctrico con pérdidas, se expresa como un número complejo:

\begin{equation}
    \varepsilon^*(\omega) = \varepsilon'(\omega) - j\,\varepsilon''(\omega),
    \label{eq:permitividad_compleja}
\end{equation}

\noindent donde:
\begin{itemize}
    \item $\varepsilon'$ es la parte real, asociada a la energía almacenada en el campo
          eléctrico (capacitiva). Cuantifica cuánto puede polarizarse el material.
    \item $\varepsilon''$ es la parte imaginaria, siempre no negativa para un medio
          pasivo, asociada a las pérdidas dieléctricas (conversión de energía
          electromagnética en calor).
    \item $\omega = 2\pi f$ es la frecuencia angular.
\end{itemize}

La permitividad relativa $\varepsilon_r$ se define normalizando respecto a la
permitividad del vacío $\varepsilon_0 = 8{,}854 \times 10^{-12}$~F/m:

\begin{equation}
    \varepsilon_r^* = \frac{\varepsilon^*}{\varepsilon_0}
                    = \varepsilon_r' - j\,\varepsilon_r''.
\end{equation}

Una magnitud derivada de uso frecuente es el \emph{factor de disipación} o
\emph{tangente de pérdidas}:

\begin{equation}
    \tan\delta = \frac{\varepsilon_r''}{\varepsilon_r'}.
\end{equation}

A lo largo de este documento se omite el subíndice $r$ salvo aclaración explícita;
todas las magnitudes son relativas.

\newpage
\Encabezado

% -----------------------------------------------------------------------------
\subsection{Modelo de Debye para Líquidos Polares}
\label{sec:debye}
% -----------------------------------------------------------------------------

El modelo de Debye~\cite{debye1929} describe la dependencia en frecuencia de la
permitividad de un dieléctrico polar cuyas moléculas poseen un único tiempo de
relajación $\tau$. Físicamente, $\tau$ cuantifica el tiempo característico de
reorientación dipolar ante un campo oscilante.

\subsubsection{Modelo de Debye simple}

\begin{equation}
    \varepsilon^*(f) = \varepsilon_\infty
                     + \frac{\varepsilon_s - \varepsilon_\infty}{1 + j\,f / f_r},
    \label{eq:debye_simple}
\end{equation}

\noindent donde:
\begin{itemize}
    \item $\varepsilon_s$ es la permitividad estática (límite $f \to 0$).
    \item $\varepsilon_\infty$ es la permitividad a alta frecuencia (límite $f \to \infty$),
          debida a polarizaciones electrónicas e iónicas.
    \item $f_r = 1/(2\pi\tau)$ es la frecuencia de relajación.
\end{itemize}

Este modelo se aplica al agua y al metanol en el presente sistema.
Los coeficientes $\{\varepsilon_s, \varepsilon_\infty, \tau\}$ dependen de la
temperatura y se tabulan a partir de los datos del NPL Report MAT~23
\cite{gregory2012} mediante ajuste polinomial cuadrático en temperatura.

\subsubsection{Modelo de Debye doble}

Algunos líquidos exhiben dos mecanismos de relajación dominantes a frecuencias
distintas (por ejemplo, el isopropanol / IPA). Se modela con dos términos de Debye
superpuestos:

\begin{equation}
    \varepsilon^*(f) = \varepsilon_\infty
                     + \frac{\varepsilon_s - \varepsilon_h}{1 + j\,f / f_{r1}}
                     + \frac{\varepsilon_h - \varepsilon_\infty}{1 + j\,f / f_{r2}},
    \label{eq:debye_doble}
\end{equation}

\noindent donde $\varepsilon_h$ es una permitividad intermedia y $f_{r1} < f_{r2}$
son las dos frecuencias de relajación.

\subsubsection{Modelo de Debye-Gamma}

Para el etanol, el NPL Report MAT~23 propone un modelo de Debye simple con
corrección lineal en frecuencia (término~$\gamma$) que aproxima una segunda
relajación en la parte superior de la banda:

\begin{equation}
    \varepsilon^*(f) = \varepsilon_h
                     + \frac{\varepsilon_s - \varepsilon_h}{1 + j\,f / f_r}
                     - j\,\gamma \left(\frac{f}{\text{GHz}}\right),
    \label{eq:debye_gamma}
\end{equation}

siendo $\gamma$ un coeficiente dimensional en GHz$^{-1}$.

En la Tabla~\ref{tab:modelos_debye} se resume el modelo utilizado para cada
líquido disponible en el sistema.

\begin{table}[h!]
\centering
\caption{Modelos de Debye por líquido de referencia implementados en el sistema.}
\label{tab:modelos_debye}
\begin{tabular}{@{}llcc@{}}
\toprule
\textbf{Líquido} & \textbf{Modelo} & \textbf{Rango T (°C)} & \textbf{Fuente} \\
\midrule
Agua (H$_2$O) & Debye simple       & 0 -- 60  & NPL MAT 23 \cite{gregory2012} \\
IPA (isopropanol) & Debye doble    & 10 -- 30 & NPL MAT 23 \cite{gregory2012} \\
Etanol         & Debye-Gamma        & 10 -- 50 & NPL MAT 23 \cite{gregory2012} \\
Metanol        & Debye simple       & 10 -- 50 & NPL MAT 23 \cite{gregory2012} \\
\bottomrule
\end{tabular}
\end{table}

\newpage
\Encabezado

% -----------------------------------------------------------------------------
\subsection{Método de la Sonda Coaxial Abierta}
\label{sec:sonda_coaxial}
% -----------------------------------------------------------------------------

\subsubsection{Principio físico}

La técnica de sonda coaxial abierta es un método de reflexión no resonante
y no destructivo para la medición de permitividad~\cite{stuchly1980, misra1990}.
Consiste en apoyar o sumergir el extremo abierto de una línea coaxial sobre
el material bajo prueba (MBP) y medir el coeficiente de reflexión~$\Gamma$
(parámetro~$S_{11}$) en un rango de frecuencias.

Cuando el extremo de la línea coaxial se abre a un semiespacio relleno con
el material, la terminación coaxial abierta actúa como una antena de apertura
en campo cercano. La admitancia en la terminación depende fuertemente de las
propiedades dieléctricas del material en contacto:

\begin{equation}
    Y(\varepsilon^*, \omega)
    = j\omega C_0\,\varepsilon^*
    + G_0\,(\varepsilon^*)^{5/2},
    \label{eq:admitancia_modelo}
\end{equation}

\noindent donde:
\begin{itemize}
    \item $C_0$ [F] es la capacitancia de la sonda en el vacío, determinada por la
          geometría del coaxial. Domina el comportamiento a bajas frecuencias.
    \item $G_0$ [S] es el coeficiente de radiación. El término $G_0\,\varepsilon^{5/2}$
          representa las pérdidas por radiación en el campo lejano; domina a altas
          frecuencias y para sondas de mayor apertura (sonda de 21~mm respecto
          a la de 3~mm).
\end{itemize}

La relación entre la admitancia y el coeficiente de reflexión medido es la
transformada de Möbius estándar de la teoría de redes:

\begin{equation}
    \Gamma = \frac{Z_0\,Y - 1}{Z_0\,Y + 1}
    \quad \Longleftrightarrow \quad
    Y = \frac{1}{Z_0}\cdot\frac{1 - \Gamma}{1 + \Gamma},
\end{equation}

\noindent con $Z_0 = 50~\Omega$.

\subsubsection{Historia y contexto del método}

El método de sonda coaxial abierta fue formalizado matemáticamente en la
década de 1980 por Stuchly et al.~\cite{stuchly1980} y Misra~\cite{misra1990},
quienes propusieron el modelo de admitancia de~\eqref{eq:admitancia_modelo}
y demostraron que el exponente $5/2$ emerge del desarrollo en serie del campo
radiado por una apertura circular en un plano conductor infinito.

La calibración con líquidos de referencia, en lugar de estándares mecánicos de
reflexión, fue introducida para eliminar la dependencia del resultado respecto
a la calidad del contacto sonda-muestra y respecto a los errores sistemáticos
del analizador de redes. La versión de dos líquidos que se documenta aquí es
consistente con la referencia UTN-FRBA 2024~\cite{utnfrba2024} en la que se basa
el presente trabajo.

\newpage
\Encabezado

% -----------------------------------------------------------------------------
\subsection{Calibración del Sistema}
\label{sec:calibracion}
% -----------------------------------------------------------------------------

El sistema presenta dos métodos de calibración con distinto número de estándares,
distintos supuestos y distinta complejidad algorítmica.

\subsubsection{Método Completo (4 estándares)}
\label{sec:metodo_completo}

Este método es denominado en la literatura como calibración de 4 estándares
del tipo SOL + 2 líquidos (Open -- Short -- Liquid\,1 -- Liquid\,2).
Los cuatro estándares son:

\begin{enumerate}
    \item \textbf{Circuito abierto (Open)}: la sonda apunta al aire.
          $\varepsilon_{\text{open}} = \varepsilon_{\text{aire}} \approx 1$.
    \item \textbf{Cortocircuito (Short)}: la apertura de la sonda se cubre con
          una lámina metálica conductora.
    \item \textbf{Líquido de referencia~1 (Ref\,1)}: por defecto IPA (isopropanol),
          de permitividad baja-moderada.
    \item \textbf{Líquido de referencia~2 (Ref\,2)}: por defecto agua destilada,
          de permitividad alta.
\end{enumerate}

Las cuatro mediciones $S_{11}$ de calibración y las permitividades conocidas
$\varepsilon_{\text{Ref1}}$, $\varepsilon_{\text{Ref2}}$ (calculadas mediante
los modelos de Debye de la Sección~\ref{sec:debye} a la temperatura de medición)
permiten determinar los parámetros del sistema $\{C_0, G_0\}$ y los tres
coeficientes de error de la caja de errores del analizador (directividad,
adaptación de fuente, seguimiento de reflexión).

La calibración completa absorbe todos los errores sistemáticos del analizador y
del cable de la sonda simultáneamente con los parámetros de la sonda misma.
No se requiere una calibración OSM previa del analizador.

\subsubsection{Método Simplificado (3 estándares)}
\label{sec:metodo_simplificado}

El método simplificado (denominado también \emph{incompleto} en la literatura)
utiliza únicamente tres estándares:

\begin{enumerate}
    \item \textbf{Open}: ídem al método completo.
    \item \textbf{Short}: ídem al método completo.
    \item \textbf{Líquido de referencia único (Ref\,1)}: por defecto agua destilada.
\end{enumerate}

La simplificación consiste en despreciar el término de radiación ($G_0 = 0$),
reduciendo el modelo de admitancia a uno puramente capacitivo:

\begin{equation}
    Y \approx j\omega C_0\,\varepsilon^*.
\end{equation}

Con esta hipótesis, la relación entre $S_{11}$ y $\varepsilon^*$ se vuelve una
transformada de Möbius (bilineal) exacta. Tres puntos conocidos en el plano~$S_{11}$
(aire, cortocircuito, líquido de referencia) determinan unívocamente la transformada,
y la extracción de $\varepsilon^*$ del material desconocido se reduce a evaluar
una fórmula cerrada (véase Sección~\ref{sec:extraccion_simplificado}).

El método simplificado es menos preciso a alta frecuencia y para sondas de mayor
apertura, pero es adecuado para exploraciones rápidas o cuando solo se dispone de
un líquido de referencia de composición conocida.

\newpage
\Encabezado

% -----------------------------------------------------------------------------
\subsection{Los Líquidos de Referencia: Selección y Criterios}
\label{sec:liquidos_referencia}
% -----------------------------------------------------------------------------

\subsubsection{Función de los líquidos de referencia}

Los líquidos de referencia cumplen una doble función:
\begin{itemize}
    \item \textbf{Calibración del analizador}: sus mediciones de $S_{11}$
          proveen puntos de la relación $\Gamma \leftrightarrow \varepsilon^*$
          con permitividades \emph{conocidas a priori} mediante los modelos
          termodinámicos de Debye.
    \item \textbf{Determinación de los parámetros de la sonda}: al conocer
          $\varepsilon^*$ se puede resolver el sistema de ecuaciones que da
          $C_0$ y $G_0$ sin necesidad de desmontarla ni medirla
          geométricamente.
\end{itemize}

\subsubsection{Por qué se necesitan dos líquidos en el método completo}

El modelo de admitancia~\eqref{eq:admitancia_modelo} tiene dos incógnitas
reales de geometría ($C_0$, $G_0$) más los tres errores del analizador,
totalizando cinco parámetros. Los estándares Open y Short aportan dos
ecuaciones, y cada líquido aporta una ecuación adicional (una por frecuencia,
pero los parámetros son constantes en frecuencia). Con un solo líquido se
tienen cuatro ecuaciones y cinco incógnitas: el sistema está indeterminado.
Con dos líquidos se completan las cinco ecuaciones necesarias.

Desde otro punto de vista: el primer líquido (Ref\,1 = IPA) fija principalmente
el coeficiente capacitivo~$C_0$, que domina la admitancia a baja frecuencia.
El segundo líquido (Ref\,2 = agua) —de permitividad considerablemente mayor—
fija principalmente el coeficiente de radiación~$G_0$, cuyo efecto resulta
apreciable a alta frecuencia y con materiales de alta permitividad.

\subsubsection{Por qué los líquidos deben ser diferentes entre sí}

Si se utilizaran dos líquidos de permitividad idéntica o muy similar,
sus mediciones de $S_{11}$ quedarían muy próximas en el plano complejo
y el sistema de ecuaciones resultaría mal condicionado numéricamente
(matriz casi singular). Esto se manifiesta como inestabilidad y ruido
amplificado en los valores de $C_0$ y $G_0$ extraídos.

Para un buen condicionamiento, los líquidos de referencia deben:
\begin{itemize}
    \item Tener partes reales $\varepsilon'$ significativamente distintas,
          idealmente separadas por al menos una década (por ejemplo,
          $\varepsilon'_{\text{IPA}} \approx 4$--$20$ frente a
          $\varepsilon'_{\text{agua}} \approx 40$--$80$ en la banda de 10~MHz a 3~GHz).
    \item Cubrir el rango de permitividades del material bajo prueba desde ambos lados
          siempre que sea posible.
    \item Tener permitividades tabuladas con alta confianza en el rango de temperatura
          y frecuencia de operación.
\end{itemize}

\subsubsection{Criterios de selección práctica}

Los líquidos de referencia ideales deben además satisfacer:
\begin{itemize}
    \item \textbf{Compatibilidad química}: no deben atacar el conector ni la
          brida de la sonda (acero inoxidable, PTFE, latón).
    \item \textbf{Bajo índice de evaporación}: la medición de $S_{11}$ tarda
          varios segundos; un líquido muy volátil puede cambiar su concentración
          durante la medición.
    \item \textbf{Homogeneidad}: deben ser líquidos puros, no mezclas cuya
          concentración sea difícil de controlar.
    \item \textbf{Datos de referencia disponibles}: se requieren coeficientes de Debye
          tabulados en función de temperatura, idealmente con trazabilidad a
          laboratorios nacionales de metrología (NPL, NIST, PTB).
\end{itemize}

En la Tabla~\ref{tab:liquidos_disponibles} se resumen las propiedades relevantes
de los cuatro líquidos disponibles en el sistema a 25~°C y 1~GHz (valores
aproximados; los valores exactos se calculan a partir de los coeficientes
tabulados).

\begin{table}[h!]
\centering
\caption{Propiedades dieléctricas aproximadas de los líquidos de referencia del sistema a 25~°C y 1~GHz.}
\label{tab:liquidos_disponibles}
\begin{tabular}{@{}lcccc@{}}
\toprule
\textbf{Líquido} & $\varepsilon'$ & $\varepsilon''$ & $\tan\delta$ & \textbf{Rol típico} \\
\midrule
Agua (H$_2$O)         & $\approx 78$  & $\approx 10$  & $\approx 0{,}13$ & Ref\,2 (alto $\varepsilon$) \\
Metanol               & $\approx 25$  & $\approx 17$  & $\approx 0{,}68$ & Ref\,2 alternativo \\
IPA (isopropanol)     & $\approx 15$  & $\approx 7$   & $\approx 0{,}47$ & Ref\,1 (bajo $\varepsilon$) \\
Etanol                & $\approx 22$  & $\approx 15$  & $\approx 0{,}68$ & Ref\,1 alternativo \\
\bottomrule
\end{tabular}
\end{table}

\noindent El par predeterminado del sistema es \textbf{IPA} (Ref\,1) + \textbf{agua} (Ref\,2),
combinación validada en la referencia de base~\cite{utnfrba2024} y ampliamente
utilizada en la literatura~\cite{stuchly1980, misra1990, gregory2012}.

\newpage
\Encabezado

% -----------------------------------------------------------------------------
\subsection{Algoritmo de Extracción de Permitividad}
\label{sec:extraccion}
% -----------------------------------------------------------------------------

\subsubsection{Método Completo: Polinomio de Grado 5}
\label{sec:extraccion_completo}

Una vez completada la calibración, los parámetros del sistema quedan condensados en
la \emph{constante de patrón} $G_n$ y en las admitancias normalizadas de los
líquidos de referencia~$y_1$, $y_2$.

\paragraph{Cálculo de $G_n$}

A partir de los cuatro $S_{11}$ de calibración y de las permitividades conocidas
$\varepsilon_1 = \varepsilon_{\text{Ref1}}$ y $\varepsilon_2 = \varepsilon_{\text{Ref2}}$,
se calculan los numeradores y denominadores:

\begin{align}
    \mathcal{A} &= (S_O - S_K)(S_2 - S_1), \nonumber \\
    \mathcal{B} &= (S_O - S_1)(S_K - S_2)\,\varepsilon_2, \nonumber \\
    \mathcal{C} &= (S_O - S_2)(S_1 - S_K)\,\varepsilon_1, \nonumber \\[4pt]
    \mathcal{D} &= (S_O - S_K)(S_2 - S_1), \nonumber \\
    \mathcal{E} &= (S_O - S_1)(S_K - S_2)\,\varepsilon_2^{5/2}, \nonumber \\
    \mathcal{F} &= (S_O - S_2)(S_1 - S_K)\,\varepsilon_1^{5/2}, \nonumber
\end{align}

\begin{equation}
    G_n = -\frac{\mathcal{A} + \mathcal{B} + \mathcal{C}}
                {\mathcal{D} + \mathcal{E} + \mathcal{F}},
    \label{eq:gn}
\end{equation}

\noindent donde $S_O$, $S_K$, $S_1$, $S_2$ son los coeficientes de reflexión
medidos para Open, Short, Ref\,1 y Ref\,2, respectivamente (todos complejos,
dependientes de la frecuencia).

Las admitancias normalizadas de los líquidos son:
\begin{align}
    y_1 &= \varepsilon_1 + G_n\,\varepsilon_1^{5/2}, \label{eq:y1} \\
    y_2 &= \varepsilon_2 + G_n\,\varepsilon_2^{5/2}. \label{eq:y2}
\end{align}

\paragraph{Término cruzado del material bajo prueba}

Dado el $S_{11}$ medido del material desconocido ($S_m$), se calcula un
\emph{término cruzado}~$T$ mediante la razón cruzada de los cuatro puntos
$\{S_m, S_K, S_1, S_2\}$:

\begin{equation}
    T = \frac{(S_m - S_1)(S_K - S_2)}{(S_m - S_K)(S_2 - S_1)}\,y_2
      + \frac{(S_m - S_2)(S_1 - S_K)}{(S_m - S_K)(S_2 - S_1)}\,y_1.
    \label{eq:third_term}
\end{equation}

\paragraph{Formulación del polinomio de grado 5}

Sustituyendo el modelo de admitancia~\eqref{eq:admitancia_modelo} normalizado
en la razón cruzada, y realizando el cambio de variable $x = \sqrt{\varepsilon_m^*}$,
se obtiene un polinomio de grado 5 en $x$:

\begin{equation}
    G_n\,x^5 + x^2 + T = 0,
    \label{eq:polinomio}
\end{equation}

que se resuelve numéricamente para cada punto en frecuencia, obteniendo cinco
raíces complejas. La permitividad se recupera como $\varepsilon_m^* = x^2$.

\paragraph{Selección de la raíz física}

De las cinco raíces del polinomio~\eqref{eq:polinomio}, solo aquellas con
\begin{equation}
    \operatorname{Re}(\varepsilon_m^*) > 0 \quad \text{y} \quad
    \operatorname{Im}(\varepsilon_m^*) \leq 0
    \label{eq:criterio_fisico}
\end{equation}
son físicamente admisibles (medio dieléctrico pasivo con la convención de
signo~\eqref{eq:permitividad_compleja}).

Cuando la condición~\eqref{eq:criterio_fisico} no resulta suficiente para
seleccionar una raíz unívoca, se aplica un algoritmo de \emph{seguimiento de
rama en frecuencia}: comenzando en la frecuencia más alta (donde los efectos
de radiación son máximos y la raíz física es más distinguible), se sigue la
rama de raíces de mayor continuidad hacia frecuencias decrecientes, apoyándose
en una extrapolación polinómica de orden~2 con ventana de 5 puntos.

Como referencia inicial se utiliza el resultado del método simplificado
(Sección~\ref{sec:extraccion_simplificado}), que proporciona una solución
aproximada sin ambigüedad de raíces.

\paragraph{Verificación de consistencia}

Al finalizar la extracción, el sistema compara el resultado del método completo
con el del simplificado en la mitad inferior de la banda. Si la divergencia
mediana supera el 25\,\%, se emite una advertencia al operador indicando posible
mal acondicionamiento de la calibración o elección inadecuada de líquidos.

\subsubsection{Método Simplificado: Transformada Bilineal}
\label{sec:extraccion_simplificado}

Con $G_n = 0$ (sin término de radiación), el modelo de admitancia se reduce a
$Y = j\omega C_0\,\varepsilon^*$, y la relación entre $S_{11}$ y $\varepsilon^*$ es
una transformada de Möbius exacta. Cuatro puntos en el plano~$S_{11}$ y sus
correspondientes en el plano~$\varepsilon$ determinan la razón cruzada.

Identificando Open ($\varepsilon_{\text{aire}} \approx 1$), Short
($\varepsilon_{\text{cortocircuito}} \to \infty$, contribución límite) y
Ref\,1 ($\varepsilon_{\text{Ref}}$), la permitividad del material desconocido
se calcula directamente como:

\begin{equation}
    \varepsilon_m^*
    = -\,c_{\text{Ref}}\,\varepsilon_{\text{Ref}}^*
      -\,c_{\text{aire}}\,\varepsilon_{\text{aire}},
    \label{eq:formula_simplificada}
\end{equation}

\noindent con coeficientes bariocéntricos:

\begin{align}
    c_{\text{Ref}}  &= \frac{(S_m - S_O)(S_K - S_{\text{Ref}})}{(S_m - S_K)(S_{\text{Ref}} - S_O)},
    \label{eq:c_ref} \\[4pt]
    c_{\text{aire}} &= \frac{(S_m - S_{\text{Ref}})(S_O - S_K)}{(S_m - S_K)(S_{\text{Ref}} - S_O)}.
    \label{eq:c_aire}
\end{align}

La identidad $c_{\text{Ref}} + c_{\text{aire}} = -1$ es una consecuencia
directa de la razón cruzada (identidad de Ptolemy), y sirve como verificación
numérica de la coherencia del cálculo.

Nótese que~\eqref{eq:formula_simplificada} es algebraicamente cerrada: no
requiere resolver polinomios ni realizar seguimiento de ramas, lo que la hace
computacionalmente trivial y robusta frente a señales ruidosas.

\newpage
\Encabezado

% -----------------------------------------------------------------------------
\subsection{Flujo Completo del Algoritmo}
\label{sec:flujo}
% -----------------------------------------------------------------------------

La Figura~\ref{fig:flujo_algoritmo} (pendiente de inclusión) resume el flujo
de datos del sistema. En forma textual, las etapas son:

\begin{enumerate}
    \item \textbf{Configuración}: selección de sonda ($C_0$ tabulated por modelo),
          temperatura de medición, líquidos de referencia y método
          (completo o simplificado).

    \item \textbf{Fase de calibración}:
    \begin{itemize}
        \item Medir o cargar los valores presetados de $S_{11}$ para Open, Short,
              Ref\,1 (y Ref\,2 en método completo).
        \item Evaluar $\varepsilon_{\text{Ref}}^*(f, T)$ mediante los modelos de
              Debye de la Sección~\ref{sec:debye} interpolando en temperatura.
        \item Calcular $G_n(f)$, $y_1(f)$, $y_2(f)$
              según~\eqref{eq:gn}--\eqref{eq:y2} (solo método completo).
    \end{itemize}

    \item \textbf{Medición del material bajo prueba}: barrido en frecuencia con
          la NanoVNA para obtener $S_{11,m}(f)$.

    \item \textbf{Extracción de $\varepsilon^*(f)$}:
    \begin{itemize}
        \item \emph{Método completo}: calcular $T(f)$ mediante~\eqref{eq:third_term},
              resolver el polinomio~\eqref{eq:polinomio} en cada frecuencia,
              seleccionar la raíz física con seguimiento de rama.
        \item \emph{Método simplificado}: evaluar directamente~\eqref{eq:formula_simplificada}.
    \end{itemize}

    \item \textbf{Salida}: curvas $\varepsilon'(f)$ y $\varepsilon''(f)$
          presentadas en la pantalla de resultados. El resultado puede exportarse
          como archivo Touchstone~`.s1p` para análisis externo.
\end{enumerate}

\newpage
\Encabezado

% -----------------------------------------------------------------------------
\subsection{Resumen Comparativo de los Métodos}
\label{sec:resumen}
% -----------------------------------------------------------------------------

\begin{table}[h!]
\centering
\caption{Comparación entre el método completo y el método simplificado.}
\label{tab:comparacion_metodos}
\begin{tabular}{@{}lll@{}}
\toprule
\textbf{Característica}         & \textbf{Método completo}              & \textbf{Método simplificado} \\
\midrule
Estándares de calibración       & 4 (Open, Short, Ref1, Ref2)           & 3 (Open, Short, Ref) \\
Líquidos de referencia          & 2 (IPA + Agua)                        & 1 (Agua) \\
Modelo de admitancia            & $j\omega C_0\varepsilon + G_0\varepsilon^{5/2}$ & $j\omega C_0\varepsilon$ \\
Extracción                      & Polinomio grado 5 + seguimiento       & Fórmula cerrada bilineal \\
Precisión a alta frecuencia     & Alta                                  & Moderada \\
Precisión para sonda de 21\,mm  & Alta                                  & Reducida \\
Complejidad computacional       & Media                                 & Baja \\
Velocidad de ensayo             & Normal                                & Rápida (exploración) \\
\bottomrule
\end{tabular}
\end{table}

% -----------------------------------------------------------------------------
%  REFERENCIAS
% -----------------------------------------------------------------------------
% (Agregar al archivo .bib del informe o usar directamente \bibitem)
%
% \begin{thebibliography}{9}
%
% \bibitem{debye1929}
% P.~Debye,
% \emph{Polar Molecules}.
% Chemical Catalog Company, New York, 1929.
%
% \bibitem{stuchly1980}
% S.~S.~Stuchly y M.~A.~Stuchly,
% ``Coaxial line reflection methods for measuring dielectric properties of biological
%   substances at radio and microwave frequencies — A review,''
% \emph{IEEE Trans. Instrum. Meas.}, vol.~29, no.~3, pp.~176--183, Sep. 1980.
%
% \bibitem{misra1990}
% D.~Misra,
% ``A quasi-static analysis of open-ended coaxial lines,''
% \emph{IEEE Trans. Microw. Theory Techn.}, vol.~35, no.~10, pp.~925--928, Oct. 1987.
%
% \bibitem{gregory2012}
% A.~P.~Gregory y R.~N.~Clarke,
% \emph{A review of RF and microwave techniques for dielectric measurements on polar liquids},
% NPL Report MAT 23. National Physical Laboratory, Teddington, UK, 2012.
%
% \bibitem{utnfrba2024}
% [COMPLETAR: autores],
% ``Mediciones de permitividad mediante una sonda coaxial,''
% UTN-FRBA, 2024. (Informe interno de referencia del proyecto.)
%
% \end{thebibliography}

%    \end{document}
% =============================================================================

\section{Marco Teórico: Caracterización Dieléctrica mediante Sonda Coaxial Abierta}

El presente capítulo desarrolla los fundamentos físicos, matemáticos y metrológicos
del método de sonda coaxial abierta empleado por el sistema NanoVNA-UTN-Toolkit
para la determinación de la permitividad eléctrica compleja de materiales líquidos.
Se presentan los dos algoritmos implementados —el método completo y el simplificado—,
los modelos de Debye utilizados para los líquidos de referencia, y los criterios de
selección de dichos líquidos.

\newpage
\Encabezado

% -----------------------------------------------------------------------------
\subsection{La Permitividad Eléctrica Compleja}
% -----------------------------------------------------------------------------

La permitividad eléctrica $\varepsilon$ es la magnitud que describe la respuesta de
un material a la aplicación de un campo eléctrico externo. En el dominio frecuencial,
para un material dieléctrico con pérdidas, se expresa como un número complejo:

\begin{equation}
    \varepsilon^*(\omega) = \varepsilon'(\omega) - j\,\varepsilon''(\omega),
    \label{eq:permitividad_compleja}
\end{equation}

\noindent donde:
\begin{itemize}
    \item $\varepsilon'$ es la parte real, asociada a la energía almacenada en el campo
          eléctrico (capacitiva). Cuantifica cuánto puede polarizarse el material.
    \item $\varepsilon''$ es la parte imaginaria, siempre no negativa para un medio
          pasivo, asociada a las pérdidas dieléctricas (conversión de energía
          electromagnética en calor).
    \item $\omega = 2\pi f$ es la frecuencia angular.
\end{itemize}

La permitividad relativa $\varepsilon_r$ se define normalizando respecto a la
permitividad del vacío $\varepsilon_0 = 8{,}854 \times 10^{-12}$~F/m:

\begin{equation}
    \varepsilon_r^* = \frac{\varepsilon^*}{\varepsilon_0}
                    = \varepsilon_r' - j\,\varepsilon_r''.
\end{equation}

Una magnitud derivada de uso frecuente es el \emph{factor de disipación} o
\emph{tangente de pérdidas}:

\begin{equation}
    \tan\delta = \frac{\varepsilon_r''}{\varepsilon_r'}.
\end{equation}

A lo largo de este documento se omite el subíndice $r$ salvo aclaración explícita;
todas las magnitudes son relativas.

\newpage
\Encabezado

% -----------------------------------------------------------------------------
\subsection{Modelo de Debye para Líquidos Polares}
\label{sec:debye}
% -----------------------------------------------------------------------------

El modelo de Debye~\cite{debye1929} describe la dependencia en frecuencia de la
permitividad de un dieléctrico polar cuyas moléculas poseen un único tiempo de
relajación $\tau$. Físicamente, $\tau$ cuantifica el tiempo característico de
reorientación dipolar ante un campo oscilante.

\subsubsection{Modelo de Debye simple}

\begin{equation}
    \varepsilon^*(f) = \varepsilon_\infty
                     + \frac{\varepsilon_s - \varepsilon_\infty}{1 + j\,f / f_r},
    \label{eq:debye_simple}
\end{equation}

\noindent donde:
\begin{itemize}
    \item $\varepsilon_s$ es la permitividad estática (límite $f \to 0$).
    \item $\varepsilon_\infty$ es la permitividad a alta frecuencia (límite $f \to \infty$),
          debida a polarizaciones electrónicas e iónicas.
    \item $f_r = 1/(2\pi\tau)$ es la frecuencia de relajación.
\end{itemize}

Este modelo se aplica al agua y al metanol en el presente sistema.
Los coeficientes $\{\varepsilon_s, \varepsilon_\infty, \tau\}$ dependen de la
temperatura y se tabulan a partir de los datos del NPL Report MAT~23
\cite{gregory2012} mediante ajuste polinomial cuadrático en temperatura.

\subsubsection{Modelo de Debye doble}

Algunos líquidos exhiben dos mecanismos de relajación dominantes a frecuencias
distintas (por ejemplo, el isopropanol / IPA). Se modela con dos términos de Debye
superpuestos:

\begin{equation}
    \varepsilon^*(f) = \varepsilon_\infty
                     + \frac{\varepsilon_s - \varepsilon_h}{1 + j\,f / f_{r1}}
                     + \frac{\varepsilon_h - \varepsilon_\infty}{1 + j\,f / f_{r2}},
    \label{eq:debye_doble}
\end{equation}

\noindent donde $\varepsilon_h$ es una permitividad intermedia y $f_{r1} < f_{r2}$
son las dos frecuencias de relajación.

\subsubsection{Modelo de Debye-Gamma}

Para el etanol, el NPL Report MAT~23 propone un modelo de Debye simple con
corrección lineal en frecuencia (término~$\gamma$) que aproxima una segunda
relajación en la parte superior de la banda:

\begin{equation}
    \varepsilon^*(f) = \varepsilon_h
                     + \frac{\varepsilon_s - \varepsilon_h}{1 + j\,f / f_r}
                     - j\,\gamma \left(\frac{f}{\text{GHz}}\right),
    \label{eq:debye_gamma}
\end{equation}

siendo $\gamma$ un coeficiente dimensional en GHz$^{-1}$.

En la Tabla~\ref{tab:modelos_debye} se resume el modelo utilizado para cada
líquido disponible en el sistema.

\begin{table}[h!]
\centering
\caption{Modelos de Debye por líquido de referencia implementados en el sistema.}
\label{tab:modelos_debye}
\begin{tabular}{@{}llcc@{}}
\toprule
\textbf{Líquido} & \textbf{Modelo} & \textbf{Rango T (°C)} & \textbf{Fuente} \\
\midrule
Agua (H$_2$O) & Debye simple       & 0 -- 60  & NPL MAT 23 \cite{gregory2012} \\
IPA (isopropanol) & Debye doble    & 10 -- 30 & NPL MAT 23 \cite{gregory2012} \\
Etanol         & Debye-Gamma        & 10 -- 50 & NPL MAT 23 \cite{gregory2012} \\
Metanol        & Debye simple       & 10 -- 50 & NPL MAT 23 \cite{gregory2012} \\
\bottomrule
\end{tabular}
\end{table}

\newpage
\Encabezado

% -----------------------------------------------------------------------------
\subsection{Método de la Sonda Coaxial Abierta}
\label{sec:sonda_coaxial}
% -----------------------------------------------------------------------------

\subsubsection{Principio físico}

La técnica de sonda coaxial abierta es un método de reflexión no resonante
y no destructivo para la medición de permitividad~\cite{stuchly1980, misra1990}.
Consiste en apoyar o sumergir el extremo abierto de una línea coaxial sobre
el material bajo prueba (MBP) y medir el coeficiente de reflexión~$\Gamma$
(parámetro~$S_{11}$) en un rango de frecuencias.

Cuando el extremo de la línea coaxial se abre a un semiespacio relleno con
el material, la terminación coaxial abierta actúa como una antena de apertura
en campo cercano. La admitancia en la terminación depende fuertemente de las
propiedades dieléctricas del material en contacto:

\begin{equation}
    Y(\varepsilon^*, \omega)
    = j\omega C_0\,\varepsilon^*
    + G_0\,(\varepsilon^*)^{5/2},
    \label{eq:admitancia_modelo}
\end{equation}

\noindent donde:
\begin{itemize}
    \item $C_0$ [F] es la capacitancia de la sonda en el vacío, determinada por la
          geometría del coaxial. Domina el comportamiento a bajas frecuencias.
    \item $G_0$ [S] es el coeficiente de radiación. El término $G_0\,\varepsilon^{5/2}$
          representa las pérdidas por radiación en el campo lejano; domina a altas
          frecuencias y para sondas de mayor apertura (sonda de 21~mm respecto
          a la de 3~mm).
\end{itemize}

La relación entre la admitancia y el coeficiente de reflexión medido es la
transformada de Möbius estándar de la teoría de redes:

\begin{equation}
    \Gamma = \frac{Z_0\,Y - 1}{Z_0\,Y + 1}
    \quad \Longleftrightarrow \quad
    Y = \frac{1}{Z_0}\cdot\frac{1 - \Gamma}{1 + \Gamma},
\end{equation}

\noindent con $Z_0 = 50~\Omega$.

\subsubsection{Historia y contexto del método}

El método de sonda coaxial abierta fue formalizado matemáticamente en la
década de 1980 por Stuchly et al.~\cite{stuchly1980} y Misra~\cite{misra1990},
quienes propusieron el modelo de admitancia de~\eqref{eq:admitancia_modelo}
y demostraron que el exponente $5/2$ emerge del desarrollo en serie del campo
radiado por una apertura circular en un plano conductor infinito.

La calibración con líquidos de referencia, en lugar de estándares mecánicos de
reflexión, fue introducida para eliminar la dependencia del resultado respecto
a la calidad del contacto sonda-muestra y respecto a los errores sistemáticos
del analizador de redes. La versión de dos líquidos que se documenta aquí es
consistente con la referencia UTN-FRBA 2024~\cite{utnfrba2024} en la que se basa
el presente trabajo.

\newpage
\Encabezado

% -----------------------------------------------------------------------------
\subsection{Calibración del Sistema}
\label{sec:calibracion}
% -----------------------------------------------------------------------------

El sistema presenta dos métodos de calibración con distinto número de estándares,
distintos supuestos y distinta complejidad algorítmica.

\subsubsection{Método Completo (4 estándares)}
\label{sec:metodo_completo}

Este método es denominado en la literatura como calibración de 4 estándares
del tipo SOL + 2 líquidos (Open -- Short -- Liquid\,1 -- Liquid\,2).
Los cuatro estándares son:

\begin{enumerate}
    \item \textbf{Circuito abierto (Open)}: la sonda apunta al aire.
          $\varepsilon_{\text{open}} = \varepsilon_{\text{aire}} \approx 1$.
    \item \textbf{Cortocircuito (Short)}: la apertura de la sonda se cubre con
          una lámina metálica conductora.
    \item \textbf{Líquido de referencia~1 (Ref\,1)}: por defecto IPA (isopropanol),
          de permitividad baja-moderada.
    \item \textbf{Líquido de referencia~2 (Ref\,2)}: por defecto agua destilada,
          de permitividad alta.
\end{enumerate}

Las cuatro mediciones $S_{11}$ de calibración y las permitividades conocidas
$\varepsilon_{\text{Ref1}}$, $\varepsilon_{\text{Ref2}}$ (calculadas mediante
los modelos de Debye de la Sección~\ref{sec:debye} a la temperatura de medición)
permiten determinar los parámetros del sistema $\{C_0, G_0\}$ y los tres
coeficientes de error de la caja de errores del analizador (directividad,
adaptación de fuente, seguimiento de reflexión).

La calibración completa absorbe todos los errores sistemáticos del analizador y
del cable de la sonda simultáneamente con los parámetros de la sonda misma.
No se requiere una calibración OSM previa del analizador.

\subsubsection{Método Simplificado (3 estándares)}
\label{sec:metodo_simplificado}

El método simplificado (denominado también \emph{incompleto} en la literatura)
utiliza únicamente tres estándares:

\begin{enumerate}
    \item \textbf{Open}: ídem al método completo.
    \item \textbf{Short}: ídem al método completo.
    \item \textbf{Líquido de referencia único (Ref\,1)}: por defecto agua destilada.
\end{enumerate}

La simplificación consiste en despreciar el término de radiación ($G_0 = 0$),
reduciendo el modelo de admitancia a uno puramente capacitivo:

\begin{equation}
    Y \approx j\omega C_0\,\varepsilon^*.
\end{equation}

Con esta hipótesis, la relación entre $S_{11}$ y $\varepsilon^*$ se vuelve una
transformada de Möbius (bilineal) exacta. Tres puntos conocidos en el plano~$S_{11}$
(aire, cortocircuito, líquido de referencia) determinan unívocamente la transformada,
y la extracción de $\varepsilon^*$ del material desconocido se reduce a evaluar
una fórmula cerrada (véase Sección~\ref{sec:extraccion_simplificado}).

El método simplificado es menos preciso a alta frecuencia y para sondas de mayor
apertura, pero es adecuado para exploraciones rápidas o cuando solo se dispone de
un líquido de referencia de composición conocida.

\newpage
\Encabezado

% -----------------------------------------------------------------------------
\subsection{Los Líquidos de Referencia: Selección y Criterios}
\label{sec:liquidos_referencia}
% -----------------------------------------------------------------------------

\subsubsection{Función de los líquidos de referencia}

Los líquidos de referencia cumplen una doble función:
\begin{itemize}
    \item \textbf{Calibración del analizador}: sus mediciones de $S_{11}$
          proveen puntos de la relación $\Gamma \leftrightarrow \varepsilon^*$
          con permitividades \emph{conocidas a priori} mediante los modelos
          termodinámicos de Debye.
    \item \textbf{Determinación de los parámetros de la sonda}: al conocer
          $\varepsilon^*$ se puede resolver el sistema de ecuaciones que da
          $C_0$ y $G_0$ sin necesidad de desmontarla ni medirla
          geométricamente.
\end{itemize}

\subsubsection{Por qué se necesitan dos líquidos en el método completo}

El modelo de admitancia~\eqref{eq:admitancia_modelo} tiene dos incógnitas
reales de geometría ($C_0$, $G_0$) más los tres errores del analizador,
totalizando cinco parámetros. Los estándares Open y Short aportan dos
ecuaciones, y cada líquido aporta una ecuación adicional (una por frecuencia,
pero los parámetros son constantes en frecuencia). Con un solo líquido se
tienen cuatro ecuaciones y cinco incógnitas: el sistema está indeterminado.
Con dos líquidos se completan las cinco ecuaciones necesarias.

Desde otro punto de vista: el primer líquido (Ref\,1 = IPA) fija principalmente
el coeficiente capacitivo~$C_0$, que domina la admitancia a baja frecuencia.
El segundo líquido (Ref\,2 = agua) —de permitividad considerablemente mayor—
fija principalmente el coeficiente de radiación~$G_0$, cuyo efecto resulta
apreciable a alta frecuencia y con materiales de alta permitividad.

\subsubsection{Por qué los líquidos deben ser diferentes entre sí}

Si se utilizaran dos líquidos de permitividad idéntica o muy similar,
sus mediciones de $S_{11}$ quedarían muy próximas en el plano complejo
y el sistema de ecuaciones resultaría mal condicionado numéricamente
(matriz casi singular). Esto se manifiesta como inestabilidad y ruido
amplificado en los valores de $C_0$ y $G_0$ extraídos.

Para un buen condicionamiento, los líquidos de referencia deben:
\begin{itemize}
    \item Tener partes reales $\varepsilon'$ significativamente distintas,
          idealmente separadas por al menos una década (por ejemplo,
          $\varepsilon'_{\text{IPA}} \approx 4$--$20$ frente a
          $\varepsilon'_{\text{agua}} \approx 40$--$80$ en la banda de 10~MHz a 3~GHz).
    \item Cubrir el rango de permitividades del material bajo prueba desde ambos lados
          siempre que sea posible.
    \item Tener permitividades tabuladas con alta confianza en el rango de temperatura
          y frecuencia de operación.
\end{itemize}

\subsubsection{Criterios de selección práctica}

Los líquidos de referencia ideales deben además satisfacer:
\begin{itemize}
    \item \textbf{Compatibilidad química}: no deben atacar el conector ni la
          brida de la sonda (acero inoxidable, PTFE, latón).
    \item \textbf{Bajo índice de evaporación}: la medición de $S_{11}$ tarda
          varios segundos; un líquido muy volátil puede cambiar su concentración
          durante la medición.
    \item \textbf{Homogeneidad}: deben ser líquidos puros, no mezclas cuya
          concentración sea difícil de controlar.
    \item \textbf{Datos de referencia disponibles}: se requieren coeficientes de Debye
          tabulados en función de temperatura, idealmente con trazabilidad a
          laboratorios nacionales de metrología (NPL, NIST, PTB).
\end{itemize}

En la Tabla~\ref{tab:liquidos_disponibles} se resumen las propiedades relevantes
de los cuatro líquidos disponibles en el sistema a 25~°C y 1~GHz (valores
aproximados; los valores exactos se calculan a partir de los coeficientes
tabulados).

\begin{table}[h!]
\centering
\caption{Propiedades dieléctricas aproximadas de los líquidos de referencia del sistema a 25~°C y 1~GHz.}
\label{tab:liquidos_disponibles}
\begin{tabular}{@{}lcccc@{}}
\toprule
\textbf{Líquido} & $\varepsilon'$ & $\varepsilon''$ & $\tan\delta$ & \textbf{Rol típico} \\
\midrule
Agua (H$_2$O)         & $\approx 78$  & $\approx 10$  & $\approx 0{,}13$ & Ref\,2 (alto $\varepsilon$) \\
Metanol               & $\approx 25$  & $\approx 17$  & $\approx 0{,}68$ & Ref\,2 alternativo \\
IPA (isopropanol)     & $\approx 15$  & $\approx 7$   & $\approx 0{,}47$ & Ref\,1 (bajo $\varepsilon$) \\
Etanol                & $\approx 22$  & $\approx 15$  & $\approx 0{,}68$ & Ref\,1 alternativo \\
\bottomrule
\end{tabular}
\end{table}

\noindent El par predeterminado del sistema es \textbf{IPA} (Ref\,1) + \textbf{agua} (Ref\,2),
combinación validada en la referencia de base~\cite{utnfrba2024} y ampliamente
utilizada en la literatura~\cite{stuchly1980, misra1990, gregory2012}.

\newpage
\Encabezado

% -----------------------------------------------------------------------------
\subsection{Algoritmo de Extracción de Permitividad}
\label{sec:extraccion}
% -----------------------------------------------------------------------------

\subsubsection{Método Completo: Polinomio de Grado 5}
\label{sec:extraccion_completo}

Una vez completada la calibración, los parámetros del sistema quedan condensados en
la \emph{constante de patrón} $G_n$ y en las admitancias normalizadas de los
líquidos de referencia~$y_1$, $y_2$.

\paragraph{Cálculo de $G_n$}

A partir de los cuatro $S_{11}$ de calibración y de las permitividades conocidas
$\varepsilon_1 = \varepsilon_{\text{Ref1}}$ y $\varepsilon_2 = \varepsilon_{\text{Ref2}}$,
se calculan los numeradores y denominadores:

\begin{align}
    \mathcal{A} &= (S_O - S_K)(S_2 - S_1), \nonumber \\
    \mathcal{B} &= (S_O - S_1)(S_K - S_2)\,\varepsilon_2, \nonumber \\
    \mathcal{C} &= (S_O - S_2)(S_1 - S_K)\,\varepsilon_1, \nonumber \\[4pt]
    \mathcal{D} &= (S_O - S_K)(S_2 - S_1), \nonumber \\
    \mathcal{E} &= (S_O - S_1)(S_K - S_2)\,\varepsilon_2^{5/2}, \nonumber \\
    \mathcal{F} &= (S_O - S_2)(S_1 - S_K)\,\varepsilon_1^{5/2}, \nonumber
\end{align}

\begin{equation}
    G_n = -\frac{\mathcal{A} + \mathcal{B} + \mathcal{C}}
                {\mathcal{D} + \mathcal{E} + \mathcal{F}},
    \label{eq:gn}
\end{equation}

\noindent donde $S_O$, $S_K$, $S_1$, $S_2$ son los coeficientes de reflexión
medidos para Open, Short, Ref\,1 y Ref\,2, respectivamente (todos complejos,
dependientes de la frecuencia).

Las admitancias normalizadas de los líquidos son:
\begin{align}
    y_1 &= \varepsilon_1 + G_n\,\varepsilon_1^{5/2}, \label{eq:y1} \\
    y_2 &= \varepsilon_2 + G_n\,\varepsilon_2^{5/2}. \label{eq:y2}
\end{align}

\paragraph{Término cruzado del material bajo prueba}

Dado el $S_{11}$ medido del material desconocido ($S_m$), se calcula un
\emph{término cruzado}~$T$ mediante la razón cruzada de los cuatro puntos
$\{S_m, S_K, S_1, S_2\}$:

\begin{equation}
    T = \frac{(S_m - S_1)(S_K - S_2)}{(S_m - S_K)(S_2 - S_1)}\,y_2
      + \frac{(S_m - S_2)(S_1 - S_K)}{(S_m - S_K)(S_2 - S_1)}\,y_1.
    \label{eq:third_term}
\end{equation}

\paragraph{Formulación del polinomio de grado 5}

Sustituyendo el modelo de admitancia~\eqref{eq:admitancia_modelo} normalizado
en la razón cruzada, y realizando el cambio de variable $x = \sqrt{\varepsilon_m^*}$,
se obtiene un polinomio de grado 5 en $x$:

\begin{equation}
    G_n\,x^5 + x^2 + T = 0,
    \label{eq:polinomio}
\end{equation}

que se resuelve numéricamente para cada punto en frecuencia, obteniendo cinco
raíces complejas. La permitividad se recupera como $\varepsilon_m^* = x^2$.

\paragraph{Selección de la raíz física}

De las cinco raíces del polinomio~\eqref{eq:polinomio}, solo aquellas con
\begin{equation}
    \operatorname{Re}(\varepsilon_m^*) > 0 \quad \text{y} \quad
    \operatorname{Im}(\varepsilon_m^*) \leq 0
    \label{eq:criterio_fisico}
\end{equation}
son físicamente admisibles (medio dieléctrico pasivo con la convención de
signo~\eqref{eq:permitividad_compleja}).

Cuando la condición~\eqref{eq:criterio_fisico} no resulta suficiente para
seleccionar una raíz unívoca, se aplica un algoritmo de \emph{seguimiento de
rama en frecuencia}: comenzando en la frecuencia más alta (donde los efectos
de radiación son máximos y la raíz física es más distinguible), se sigue la
rama de raíces de mayor continuidad hacia frecuencias decrecientes, apoyándose
en una extrapolación polinómica de orden~2 con ventana de 5 puntos.

Como referencia inicial se utiliza el resultado del método simplificado
(Sección~\ref{sec:extraccion_simplificado}), que proporciona una solución
aproximada sin ambigüedad de raíces.

\paragraph{Verificación de consistencia}

Al finalizar la extracción, el sistema compara el resultado del método completo
con el del simplificado en la mitad inferior de la banda. Si la divergencia
mediana supera el 25\,\%, se emite una advertencia al operador indicando posible
mal acondicionamiento de la calibración o elección inadecuada de líquidos.

\subsubsection{Método Simplificado: Transformada Bilineal}
\label{sec:extraccion_simplificado}

Con $G_n = 0$ (sin término de radiación), el modelo de admitancia se reduce a
$Y = j\omega C_0\,\varepsilon^*$, y la relación entre $S_{11}$ y $\varepsilon^*$ es
una transformada de Möbius exacta. Cuatro puntos en el plano~$S_{11}$ y sus
correspondientes en el plano~$\varepsilon$ determinan la razón cruzada.

Identificando Open ($\varepsilon_{\text{aire}} \approx 1$), Short
($\varepsilon_{\text{cortocircuito}} \to \infty$, contribución límite) y
Ref\,1 ($\varepsilon_{\text{Ref}}$), la permitividad del material desconocido
se calcula directamente como:

\begin{equation}
    \varepsilon_m^*
    = -\,c_{\text{Ref}}\,\varepsilon_{\text{Ref}}^*
      -\,c_{\text{aire}}\,\varepsilon_{\text{aire}},
    \label{eq:formula_simplificada}
\end{equation}

\noindent con coeficientes bariocéntricos:

\begin{align}
    c_{\text{Ref}}  &= \frac{(S_m - S_O)(S_K - S_{\text{Ref}})}{(S_m - S_K)(S_{\text{Ref}} - S_O)},
    \label{eq:c_ref} \\[4pt]
    c_{\text{aire}} &= \frac{(S_m - S_{\text{Ref}})(S_O - S_K)}{(S_m - S_K)(S_{\text{Ref}} - S_O)}.
    \label{eq:c_aire}
\end{align}

La identidad $c_{\text{Ref}} + c_{\text{aire}} = -1$ es una consecuencia
directa de la razón cruzada (identidad de Ptolemy), y sirve como verificación
numérica de la coherencia del cálculo.

Nótese que~\eqref{eq:formula_simplificada} es algebraicamente cerrada: no
requiere resolver polinomios ni realizar seguimiento de ramas, lo que la hace
computacionalmente trivial y robusta frente a señales ruidosas.

\newpage
\Encabezado

% -----------------------------------------------------------------------------
\subsection{Flujo Completo del Algoritmo}
\label{sec:flujo}
% -----------------------------------------------------------------------------

La Figura~\ref{fig:flujo_algoritmo} (pendiente de inclusión) resume el flujo
de datos del sistema. En forma textual, las etapas son:

\begin{enumerate}
    \item \textbf{Configuración}: selección de sonda ($C_0$ tabulated por modelo),
          temperatura de medición, líquidos de referencia y método
          (completo o simplificado).

    \item \textbf{Fase de calibración}:
    \begin{itemize}
        \item Medir o cargar los valores presetados de $S_{11}$ para Open, Short,
              Ref\,1 (y Ref\,2 en método completo).
        \item Evaluar $\varepsilon_{\text{Ref}}^*(f, T)$ mediante los modelos de
              Debye de la Sección~\ref{sec:debye} interpolando en temperatura.
        \item Calcular $G_n(f)$, $y_1(f)$, $y_2(f)$
              según~\eqref{eq:gn}--\eqref{eq:y2} (solo método completo).
    \end{itemize}

    \item \textbf{Medición del material bajo prueba}: barrido en frecuencia con
          la NanoVNA para obtener $S_{11,m}(f)$.

    \item \textbf{Extracción de $\varepsilon^*(f)$}:
    \begin{itemize}
        \item \emph{Método completo}: calcular $T(f)$ mediante~\eqref{eq:third_term},
              resolver el polinomio~\eqref{eq:polinomio} en cada frecuencia,
              seleccionar la raíz física con seguimiento de rama.
        \item \emph{Método simplificado}: evaluar directamente~\eqref{eq:formula_simplificada}.
    \end{itemize}

    \item \textbf{Salida}: curvas $\varepsilon'(f)$ y $\varepsilon''(f)$
          presentadas en la pantalla de resultados. El resultado puede exportarse
          como archivo Touchstone~`.s1p` para análisis externo.
\end{enumerate}

\newpage
\Encabezado

% -----------------------------------------------------------------------------
\subsection{Resumen Comparativo de los Métodos}
\label{sec:resumen}
% -----------------------------------------------------------------------------

\begin{table}[h!]
\centering
\caption{Comparación entre el método completo y el método simplificado.}
\label{tab:comparacion_metodos}
\begin{tabular}{@{}lll@{}}
\toprule
\textbf{Característica}         & \textbf{Método completo}              & \textbf{Método simplificado} \\
\midrule
Estándares de calibración       & 4 (Open, Short, Ref1, Ref2)           & 3 (Open, Short, Ref) \\
Líquidos de referencia          & 2 (IPA + Agua)                        & 1 (Agua) \\
Modelo de admitancia            & $j\omega C_0\varepsilon + G_0\varepsilon^{5/2}$ & $j\omega C_0\varepsilon$ \\
Extracción                      & Polinomio grado 5 + seguimiento       & Fórmula cerrada bilineal \\
Precisión a alta frecuencia     & Alta                                  & Moderada \\
Precisión para sonda de 21\,mm  & Alta                                  & Reducida \\
Complejidad computacional       & Media                                 & Baja \\
Velocidad de ensayo             & Normal                                & Rápida (exploración) \\
\bottomrule
\end{tabular}
\end{table}

% -----------------------------------------------------------------------------
%  REFERENCIAS
% -----------------------------------------------------------------------------
% (Agregar al archivo .bib del informe o usar directamente \bibitem)
%
% \begin{thebibliography}{9}
%
% \bibitem{debye1929}
% P.~Debye,
% \emph{Polar Molecules}.
% Chemical Catalog Company, New York, 1929.
%
% \bibitem{stuchly1980}
% S.~S.~Stuchly y M.~A.~Stuchly,
% ``Coaxial line reflection methods for measuring dielectric properties of biological
%   substances at radio and microwave frequencies — A review,''
% \emph{IEEE Trans. Instrum. Meas.}, vol.~29, no.~3, pp.~176--183, Sep. 1980.
%
% \bibitem{misra1990}
% D.~Misra,
% ``A quasi-static analysis of open-ended coaxial lines,''
% \emph{IEEE Trans. Microw. Theory Techn.}, vol.~35, no.~10, pp.~925--928, Oct. 1987.
%
% \bibitem{gregory2012}
% A.~P.~Gregory y R.~N.~Clarke,
% \emph{A review of RF and microwave techniques for dielectric measurements on polar liquids},
% NPL Report MAT 23. National Physical Laboratory, Teddington, UK, 2012.
%
% \bibitem{utnfrba2024}
% [COMPLETAR: autores],
% ``Mediciones de permitividad mediante una sonda coaxial,''
% UTN-FRBA, 2024. (Informe interno de referencia del proyecto.)
%
% \end{thebibliography}
